\documentclass[11pt]{article}

\usepackage[a4paper,margin=1in]{geometry}
\usepackage{amsmath,amssymb,amsthm,mathtools}
\usepackage{bm}
\usepackage{enumitem}
\usepackage{url}
\usepackage{hyperref}

\newtheorem{theorem}{Theorem}[section]
\newtheorem{lemma}[theorem]{Lemma}
\newtheorem{corollary}[theorem]{Corollary}
\newtheorem{proposition}[theorem]{Proposition}
\newtheorem{definition}[theorem]{Definition}
\newtheorem{protocol}[theorem]{Protocol}
\newtheorem{remark}[theorem]{Remark}

\newcommand{\F}{\mathbb{F}}
\newcommand{\Z}{\mathbb{Z}}
\newcommand{\R}{\mathbb{F}_r}
\newcommand{\Fp}{\mathbb{F}_p}
\newcommand{\Ep}{\mathbb{F}_{p^4}}
\newcommand{\T}{\mathcal{T}}
\newcommand{\Cset}{\mathcal{C}}
\newcommand{\bits}{\{0,1\}}
\newcommand{\wt}[1]{\widehat{#1}}
\newcommand{\poly}{\mathrm{poly}}
\newcommand{\negl}{\mathrm{negl}}

\title{A Rowwise Exactness Repair for Deferred-Quotient Constraints in BN254 Wrappers}
\author{}
\date{}

\begin{document}
\maketitle

\begin{abstract}
This note studies a scope-limited class of deferred-quotient constraints in
BN254 wrappers: non-negative normalized row templates, unsigned quotient ranges,
and designated residue rows. In this scope, wrappers enforce identities of the
form \(L=D+pQ\) inside \(R=\F_r\), with \(p:=p_{\mathrm{M31}}\). Since \(p\) is
invertible in \(R\), such identities are rowwise-vacuous on \(\Omega\) unless
residue and quotient witnesses are additionally bound to small canonical integer
ranges. We present a native BN254 canonical-lift binding layer (residue,
quotient, and carry range checks plus audited no-wrap bounds) that, under
explicit rowwise premises, upgrades satisfied field equalities to exact
Euclidean-division statements in \(\Z\) on designated residue rows, thereby
excluding the rowwise \(p^{-1}\)-witness vacuity class in this restricted
model. The synthesis interface is exposed as explicit public certificates
(audit/coverage/wiring) with a partially machine-checkable artifact/checker:
the released checker verifies bounds, wiring, row-activity metadata, and digest
integrity, while template/realization coherence remains a published symbolic
proof obligation. Accordingly, the main result is a certified rowwise exactness
theorem for a restricted class of deferred-quotient BN254 wrapper
arithmetizations, not a backend-independent wrapper soundness theorem. Full
wrapper-level cryptographic closure (transcript binding, degree discipline,
Fiat--Shamir soundness, backend composition) remains backend-dependent and is
out of scope. This note targets a class of deferred-quotient wrapper
arithmetizations over BN254; it does not claim a flaw in BN254 itself or in
Halo2 as a proof system.
\end{abstract}

\section{Deferred-Quotient Vacuity and a Rowwise Exactness Repair}

Let
\[
p := p_{\mathrm{M31}} = 2^{31}-1,
\qquad
R := \F_r,
\]
where \(R\) is the BN254 scalar field. We assume the one-limb headroom condition
\[
p\cdot 2^{66}+p < r.
\]
For \(p=2^{31}-1\), the left-hand side equals
\[
(2^{31}-1)(2^{66}+1)=2^{97}+2^{31}-2^{66}-1<2^{97},
\]
so this condition holds for BN254 scalar \(r\). This class-level headroom
condition is sufficient for the one-limb \(31/66\)-bit quotient classes used
below. Full family-specific left/right no-wrap bounds are discharged later by
the public audit certificate in Appendix~\ref{app:audit-certificate}.

Let
\[
\Omega = \{\omega^0,\omega^1,\dots,\omega^{m-1}\} \subset R
\]
be the public row domain indexing the target transcript tables, with size
\[
m := |\Omega|.
\]

We keep the deferred-quotient compression layer unchanged. For each scalar local
constraint index \(t\), let
\[
L_t(X), D_t(X) \in R[X]
\]
denote the (potentially high-degree) \emph{composition polynomials} induced by
the instantiated row templates, i.e., obtained by evaluating the row-template
polynomials on the underlying degree-\(<m\) witness polynomials over \(\Omega\).
Let
\[
Q_t(X) \in R[X]
\]
denote the low-degree quotient witness polynomial (so \(\deg Q_t < m\)) that
interpolates the quotient row table on \(\Omega\). Define the aggregated
numerator
\[
N(X) := \sum_t \lambda_t\bigl(L_t(X)-D_t(X)-pQ_t(X)\bigr),
\]
where \(\lambda_t\) are random aggregation coefficients sampled by the wrapper
verifier (or via Fiat--Shamir in a non-interactive instantiation). Let
\[
Z_\Omega(X) := \prod_{u\in\Omega}(X-u)
\]
and let \(H(X)\in R[X]\) denote a prover-supplied witness polynomial. Whenever
\(N\) vanishes on \(\Omega\), define the \emph{true} quotient
\[
H_{\mathrm{true}}(X) := \frac{N(X)}{Z_\Omega(X)}.
\]

\paragraph{Composition versus interpolant.}
We emphasize that \(L_t(X)\) and \(D_t(X)\) above are \emph{not} defined as the
unique degree-\(<m\) interpolants of the rowwise values \((L_t[j])_{j\in\Omega}\)
and \((D_t[j])_{j\in\Omega}\). Under quadratic templates, the composition
polynomials can have degree up to \(2m-2\), so \(N\) is genuinely high-degree
and \(H_{\mathrm{true}}\) need not be identically zero on honest transcripts.

\paragraph{Scope and claim boundary.}
This manuscript is organized as a two-layer package:
\begin{enumerate}[label=(\roman*),leftmargin=2em]
    \item a rowwise exactness layer (canonical residues, quotient/carry range
    binding, audited no-wrap bounds) culminating in
    Theorem~\ref{thm:rowwise-no-vacuity};
    \item a brief discussion (Remark~\ref{rem:wrapper-extension}) of how the
    repair composes with a full wrapper system in concrete backends.
\end{enumerate}
The theorem proved here is a rowwise algebraic exactness theorem under explicit
public synthesis certificates (audit bounds, realization coherence, and wiring
coverage), not a full backend-independent wrapper soundness theorem.
Accordingly, this note does \emph{not} claim a fully closed stand-alone
cryptographic theorem for any particular internal-witness wrapper backend.
Closing transcript-binding, degree, and Fiat--Shamir/soundness assumptions is a
separate compilation task.

\paragraph{Restricted model covered by this note.}
The claims proved here apply only to the restricted deferred-quotient model
explicitly described in this manuscript: normalized non-negative row templates,
unsigned quotient ranges, designated residue rows, and the audited instantiated
family set recorded in the public certificate. In the current manuscript
instance, this means the audited families
\[
\{\mathrm{inv},\mathrm{EF},\mathrm{EE},\mathrm{hor},\mathrm{car}\}.
\]
Families whose satisfying assignments require signed quotient ranges, negative
normalizations, or a different residue convention are outside the scope of this
repair note. Thus the main theorem should be read as a certified
restricted-model theorem, not as a generic statement about arbitrary BN254
wrapper arithmetizations.

\subsection{Vulnerability: deferred-quotient vacuity in \(R\)}

The deferred-quotient layer aims to represent rowwise Euclidean division by
\(p\) by enforcing field equalities of the form
\[
L_t(\omega^j)-D_t(\omega^j)=pQ_t(\omega^j)
\]
for all \(\omega^j\in\Omega\), where \(Q_t(X)\) is the degree-\(<m\) quotient
witness polynomial interpolating the row-table values
\((Q_t[j])_{j\in\Omega}\).
Since \(p\) is invertible in \(R\), this rowwise constraint on \(\Omega\) is
vacuous unless quotient values are additionally range-bound: at the level of
actual composition evaluations \((L_t(\omega^j),D_t(\omega^j))_{j\in\Omega}\),
the prover can set
\[
Q_t[j]:=p^{-1}\bigl(L_t(\omega^j)-D_t(\omega^j)\bigr)\in R
\]
so that \(L_t(\omega^j)-D_t(\omega^j)=pQ_t(\omega^j)\) holds on \(\Omega\).
There is then a unique degree-\(<m\) interpolant through the chosen quotient
row values, so the rowwise constraint carries no intended integer-division
semantics unless quotient cells are bound to small integer ranges.

\begin{definition}[Compression/rowwise interface consistency]
\label{def:compression-rowwise-rhs-consistency}
For each relation index \(t\), let \((L_t[j],D_t[j],Q_t[j])_{j\in\Omega}\) be the
rowwise values defined by the realization rule of
Definition~\ref{def:selector-relation-realization} together with the
zero-extension rule of Definition~\ref{def:rowtable-zero-extension}.
The wrapper compression layer uses the composition polynomials \(L_t(X),D_t(X)\)
from the instantiated templates (as in the numerator definition above), and the
quotient witness polynomial \(Q_t(X)\) interpolating \((Q_t[j])_{j\in\Omega}\).
Consistency requires that
\[
L_t(\omega^j)=L_t[j],
\qquad
D_t(\omega^j)=D_t[j],
\qquad
Q_t(\omega^j)=Q_t[j]
\quad\text{for all }j\in\{0,\dots,|\Omega|-1\}.
\]
\end{definition}

Under Definition~\ref{def:compression-rowwise-rhs-consistency}, this rowwise
vacuity extends directly to the polynomial quotient-check layer.

\begin{proposition}[Polynomial-level vacuity under unbounded quotient witnesses]
\label{prop:poly-level-vacuity}
Fix a relation index \(t\) and a witness assignment determining composition
polynomials
\[
L_t(X),D_t(X)\in R[X]
\]
with
\[
\deg L_t,\deg D_t \le 2m-2.
\]
Assume quotient witnesses are not range-bound. For each row \(\omega^j\in\Omega\),
define
\[
Q_t[j] := p^{-1}\bigl(L_t(\omega^j)-D_t(\omega^j)\bigr)\in R,
\]
let \(Q_t(X)\) be the unique degree-\(<m\) interpolant through the row values
\((Q_t[j])_{j\in\Omega}\), and set
\[
R_t(X):=L_t(X)-D_t(X)-pQ_t(X).
\]
Then \(R_t(\omega^j)=0\) for all \(j\in\{0,\dots,m-1\}\), so \(Z_\Omega\mid R_t\).
Moreover
\[
\deg R_t\le 2m-2,
\]
hence the quotient
\[
H_t(X):=\frac{R_t(X)}{Z_\Omega(X)}
\]
has degree \(\le m-2\). Therefore the usual low-degree quotient-check layer
alone does not enforce intended Euclidean-division semantics.
\end{proposition}

\begin{proof}
For each \(j\),
\[
pQ_t(\omega^j)=pQ_t[j]=L_t(\omega^j)-D_t(\omega^j)
\]
by construction of the interpolated quotient values. Hence
\[
R_t(\omega^j)=L_t(\omega^j)-D_t(\omega^j)-pQ_t(\omega^j)=0
\]
for all \(\omega^j\in\Omega\), so \(Z_\Omega\mid R_t\). The degree bound follows
from
\[
\deg L_t,\deg D_t\le 2m-2,\qquad \deg(pQ_t)<m.
\]
Since \(\deg Z_\Omega=m\), the quotient degree bound \(\deg H_t\le m-2\) follows.
\end{proof}

\paragraph{Evaluation-level formulation.}
Proposition~\ref{prop:poly-level-vacuity} is intentionally stated directly in
terms of the actual composition evaluations
\(L_t(\omega^j),D_t(\omega^j)\). It therefore isolates the vacuity phenomenon
without depending on any later row-table serialization convention.

\begin{corollary}[Aggregated numerator vacuity under unbounded quotient witnesses]
\label{cor:aggregated-vacuity}
Let \(\mathcal I\) be a finite relation-index set. For each \(t\in\mathcal I\),
assume the setup of Proposition~\ref{prop:poly-level-vacuity}, and define
\[
R_t(X):=L_t(X)-D_t(X)-pQ_t(X).
\]
For coefficients \((\lambda_t)_{t\in\mathcal I}\subset R\), set
\[
N(X):=\sum_{t\in\mathcal I}\lambda_t R_t(X).
\]
Then \(Z_\Omega\mid N\). If additionally
\[
\deg L_t,\deg D_t\le 2m-2,\qquad \deg Q_t<m \quad (\forall t\in\mathcal I),
\]
then \(\deg N\le 2m-2\), and the quotient
\[
H_N(X):=\frac{N(X)}{Z_\Omega(X)}
\]
has degree \(\le m-2\).
\end{corollary}

\begin{proof}
For each \(t\), Proposition~\ref{prop:poly-level-vacuity} gives \(Z_\Omega\mid R_t\).
Hence \(Z_\Omega\) divides every \(R\)-linear combination \(\sum_t\lambda_tR_t\),
proving \(Z_\Omega\mid N\). The degree bound follows from
\(\deg R_t\le 2m-2\) for each \(t\), so \(\deg N\le 2m-2\), and then
\[
\deg H_N\le (2m-2)-m=m-2.
\]
\end{proof}

\paragraph{Terminology.}
We refer to this pattern as a \emph{deferred-quotient} wrapper: modular
reductions are enforced via identities of the form \(L=D+pQ\) in the outer
field, deferring explicit non-native division to a range-binding layer.

\paragraph{Toy counterexample (\(p^{-1}\) witness).}
At the row level, if the circuit only enforces \(L=D+pQ\) in \(R\) with an
unbounded quotient witness, a malicious prover can satisfy the check while
violating the intended modulo-\(p\) semantics. For instance, take \(L=0\) and
\(D=1\) in \(R\), and set \(Q:=-p^{-1}\bmod r\). Then \(0=1+p(-p^{-1})\) holds in
\(R\) even though it does not encode any valid integer division statement.

\subsection{Repair summary and claim level}

The missing ingredient is a lightweight realization in this model of the canonical-lift
property \(\mathcal{P}_{\mathrm{canon}}\): residue cells must be bound to
their canonical integer ranges, and quotient/carry cells must be bounded as
small integers, so that a field equality in \(R\) upgrades to an exact
Euclidean-division statement in \(\Z\).
\begin{definition}[Canonical-lift property guarantee]
\label{def:f-canon}
We say a binding layer realizes the canonical-lift property guarantee
\(\mathcal{P}_{\mathrm{canon}}\) if, for every satisfying assignment:
\begin{enumerate}[label=(\roman*),leftmargin=2em]
    \item every designated residue cell is the image in \(R\) of an integer in
    \(I_p=\{0,\dots,p-1\}\);
    \item every designated quotient/carry cell is the image in \(R\) of an
    integer in its declared small range (e.g., \(I_{31}\), \(I_{66}\),
    \(\{0,1\}\), or \(\{0,\dots,7\}\));
    \item for every satisfied deferred-quotient row relation
    \(L_t[j]=D_t[j]+pQ_t[j]\) in \(R\), the row values admit canonical integer
    lifts \(\wt{L_t[j]},\wt{D_t[j]},\wt{Q_t[j]}\) such that
    \[
    0\le \wt{L_t[j]}<r,
    \qquad
    0\le \wt{D_t[j]}+p\wt{Q_t[j]}<r;
    \]
    hence the relation holds as an exact identity in \(\Z\), and on designated
    residue rows it encodes Euclidean division by \(p\).
\end{enumerate}
\end{definition}
Concretely, this note realizes \(\mathcal{P}_{\mathrm{canon}}\) by adding only:
\begin{enumerate}[label=(\roman*),leftmargin=2em]
    \item exact residue canonicality in \([0,p-1]\) using \(31\)-bit limb
    decomposition plus one exclusion gate;
    \item coarse \(31\)-bit or \(66\)-bit bounds for quotient witnesses, and
    tiny bounds for carries;
    \item audited no-wrap bounds for all satisfying witnesses.
\end{enumerate}
These yield the main algebraic repair result, Theorem~\ref{thm:rowwise-no-vacuity}.
We include only a brief remark on composing the repair with a full wrapper
system (Remark~\ref{rem:wrapper-extension}).

\subsection{Positioning and related work}

This repair follows a standard non-native (foreign-field) arithmetic pattern:
express modular reduction via an equality carrying an explicit
quotient/remainder witness, and then range-bind those witnesses to small
integer intervals so the equality has a unique \(\Z\)-semantics. This pattern
is ubiquitous in PLONKish arithmetizations and their lookup-based range-check
extensions~\cite{plonk2019,plookup2020,lookup-sok-2025}.
Recent work on foreign-field arithmetic in PLONK studies tight headroom/no-wrap
premises and their constraint costs under such bounds~\cite{foreignfield2025}.
For a representative engineering codebase that includes non-native arithmetic gadgets in
Halo2 circuits, see~\cite{halo2wrong}.

Accordingly, we present this note as a simple \emph{algebraic repair} for
deferred-quotient BN254 wrappers: our novelty is not a new foreign-field gadget
per se, but the diagnosis and repair of a specific algebraic vacuity (the
\(p^{-1}\) attack) via an explicit canonical-lift binding and public audit/wiring
certificates.

\subsection{Column classes}

We partition the witness columns into five classes:
\begin{itemize}[leftmargin=2em]
    \item \(\Cset_p\): residue columns that must lie in
    \[
    I_p := \{0,1,\dots,p-1\}.
    \]
    \item \(\Cset_{31}\): quotient columns that only need the coarse bound
    \[
    I_{31} := \{0,1,\dots,2^{31}-1\}.
    \]
    \item \(\Cset_{66}\): quotient columns that only need the coarse bound
    \[
    I_{66} := \{0,1,\dots,2^{66}-1\}.
    \]
    \item \(\Cset_{\mathrm{bit}}\): carry columns in \(\{0,1\}\).
    \item \(\Cset_{8}\): small carry columns in \(\{0,1,\dots,7\}\).
\end{itemize}

The only exact canonicality requirement is for \(\Cset_p\). For quotient
columns, it is enough to enforce a public upper bound strictly below \(r/p\) so
that, together with RHS no-wrap and canonical remainder binding on designated
residue rows, any equality of the form
\[
L = D + pQ
\quad\text{in }R
\]
upgrades to an exact identity in \(\Z\). On those designated residue rows, this
exact identity yields the unique Euclidean quotient and canonical remainder.

\begin{definition}[Complete quotient coverage]
\label{def:complete-quotient-coverage}
Let \(\mathcal Q\) be the set of all quotient witness columns used by the
target deferred-quotient arithmetization. There is a public class map
\[
\kappa:\mathcal Q \to \{31,66\},
\]
with induced partition
\[
\mathcal Q_{31}:=\kappa^{-1}(31),\qquad
\mathcal Q_{66}:=\kappa^{-1}(66),
\]
such that
\[
\mathcal Q_{31}\cap\mathcal Q_{66}=\varnothing,\qquad
\mathcal Q_{31}\cup\mathcal Q_{66}=\mathcal Q.
\]
Every row-local quotient symbol \(Q_t[j]\) is wired to exactly one column in
\(\mathcal Q\). Columns in \(\mathcal Q_{31}\) are constrained by
\eqref{eq:q31-recompose}--\eqref{eq:q31-lookups}; columns in \(\mathcal Q_{66}\)
are constrained by \eqref{eq:q66-recompose}--\eqref{eq:q66-lookups}.
\end{definition}

\paragraph{Certificate metadata (moved to appendix).}
The deferred-quotient arithmetization comes with explicit public inventory and
wiring certificates (quotient-class map, selector-level slot wiring, etc.). For
readability, we move this compiler metadata to
Appendix~\ref{app:coverage-and-wiring}.

\subsection{Binding mechanism}
\label{subsec:binding-mechanism}

We now define the minimal binding constraints that realize \(\mathcal{P}_{\mathrm{canon}}\).

\subsubsection{Exact canonical residues in \([0,p-1]\)}

For every cell \(x = W[j]\) with \(W\in\Cset_p\), introduce two limb witnesses
\[
x^{(0)},x^{(1)} \in R
\]
and one exclusion witness
\[
\nu_x \in R.
\]
The circuit enforces
\begin{equation}
x = x^{(0)} + 2^{16}x^{(1)},
\label{eq:residue-recompose}
\end{equation}
\begin{equation}
x^{(0)} \in T_{16} := \{0,1,\dots,2^{16}-1\},
\qquad
x^{(1)} \in T_{15} := \{0,1,\dots,2^{15}-1\},
\label{eq:residue-limb-lookups}
\end{equation}
and
\begin{equation}
(x-p)\nu_x = 1.
\label{eq:residue-not-p}
\end{equation}

Equations \eqref{eq:residue-recompose}--\eqref{eq:residue-limb-lookups} imply
\[
0 \le x \le 2^{31}-1 = p.
\]
Equation \eqref{eq:residue-not-p} excludes the unique bad endpoint \(x=p\), so \(x\in I_p\).

\begin{lemma}[Mersenne canonicality]
\label{lem:mersenne-canonicality}
If \(x\in R\) satisfies \eqref{eq:residue-recompose}--\eqref{eq:residue-not-p}, then \(x\) is the canonical integer lift of a unique element of \(\Fp\).
\end{lemma}

\begin{proof}
From \eqref{eq:residue-recompose}--\eqref{eq:residue-limb-lookups}, we obtain \(x\in[0,2^{31}-1]\). Since \(2^{31}-1=p\), the only value in this interval that is not in \(I_p=\{0,\dots,p-1\}\) is \(x=p\). Equation \eqref{eq:residue-not-p} implies \(x\neq p\). Hence \(x\in[0,p-1]\), so it is the canonical lift of a unique residue modulo \(p\).
\end{proof}

\subsubsection{Small quotient columns}

For every cell \(q=W[j]\) with \(W\in\Cset_{31}\), introduce two limb witnesses
\[
q^{(0)},q^{(1)}\in R
\]
and enforce
\begin{equation}
q = q^{(0)} + 2^{16}q^{(1)},
\label{eq:q31-recompose}
\end{equation}
\begin{equation}
q^{(0)} \in T_{16},
\qquad
q^{(1)} \in T_{15}.
\label{eq:q31-lookups}
\end{equation}
Hence
\[
0 \le q < 2^{31}.
\]

\subsubsection{Large quotient columns}

For every cell \(q=W[j]\) with \(W\in\Cset_{66}\), introduce five limb witnesses
\[
q^{(0)},q^{(1)},q^{(2)},q^{(3)},q^{(4)} \in R
\]
and enforce
\begin{equation}
q = q^{(0)} + 2^{16}q^{(1)} + 2^{32}q^{(2)} + 2^{48}q^{(3)} + 2^{64}q^{(4)},
\label{eq:q66-recompose}
\end{equation}
\begin{equation}
q^{(0)},q^{(1)},q^{(2)},q^{(3)} \in T_{16},
\qquad
q^{(4)} \in T_{2} := \{0,1,2,3\}.
\label{eq:q66-lookups}
\end{equation}
Hence
\[
0 \le q < 2^{66}.
\]

\subsubsection{Carry columns}

For every cell \(c=W[j]\) with \(W\in\Cset_{\mathrm{bit}}\), enforce
\begin{equation}
c(c-1)=0.
\label{eq:bit-carry}
\end{equation}
For every cell \(u=W[j]\) with \(W\in\Cset_{8}\), enforce
\begin{equation}
u \in T_8 := \{0,1,\dots,7\}.
\label{eq:u8-lookup}
\end{equation}

\subsubsection{Batched lookup realization}

All instances of
\[
x^{(0)}\in T_{16},\;
x^{(1)}\in T_{15},\;
q_{31}^{(0)}\in T_{16},\;
q_{31}^{(1)}\in T_{15},\;
q_{66}^{(0)}\in T_{16},\;
q_{66}^{(1)}\in T_{16},\;
q_{66}^{(2)}\in T_{16},\;
q_{66}^{(3)}\in T_{16},\;
q^{(4)}\in T_2,\;
u\in T_8
\]
are enforced using one batched native lookup layer in the Tier-2 PLONKish circuit. Thus the binding mechanism uses only:
\begin{itemize}[leftmargin=2em]
    \item native BN254 lookups into tiny public tables \(T_{16},T_{15},T_2,T_8\),
    \item one non-equality constraint \((x-p)\nu_x=1\) per residue cell, and
    \item native linear recomposition constraints.
\end{itemize}
No foreign-field multiplication gadget, no CRT consistency check, and no multi-limb carry chain is introduced.

\subsection{Why the \(p^{-1}\) attack is eliminated}

The previous vacuity attack exploited the fact that a malicious prover could choose an unconstrained field element
\[
Q = -p^{-1}\pmod r.
\]
This is impossible once the quotient cells are bounded as small integers.

\begin{lemma}[Bounded deferred-quotient exactness]
\label{lem:bounded-quotient-exactness}
Let \(B\in\{2^{31},2^{66}\}\). Let \(L,D,q\in R\) satisfy:
\begin{enumerate}[label=(\roman*),leftmargin=2em]
    \item \(L\) is the image in \(R\) of an integer \(\wt{L}\) with
    \[
    0 \le \wt{L} < r;
    \]
    \item \(D\) is the image in \(R\) of an integer \(\wt{D}\) with
    \[
    0 \le \wt{D} < r;
    \]
    \item \(q\) is the image in \(R\) of an integer \(\wt{q}\in[0,B)\);
    \item in \(\Z\),
    \[
    0\le \wt{D}+p\wt{q}<r;
    \]
    \item in \(R\),
    \[
    L = D + pq.
    \]
\end{enumerate}
Then
\[
\wt{L}=\wt{D}+p\wt{q}
\quad\text{in }\Z.
\]
If additionally \(\wt{D}\in I_p\), then necessarily
\[
\wt{D} = \wt{L} \bmod p,
\qquad
\wt{q} = \left\lfloor \frac{\wt{L}}{p} \right\rfloor.
\]
\end{lemma}

\begin{proof}
By (i) and (iv), both \(\wt{L}\) and \(\wt{D}+p\wt{q}\) lie in \([0,r)\). Using
(v),
\[
L-(D+pq)=0 \text{ in }R
\]
upgrades to integer equality:
\[
\wt{L}=\wt{D}+p\wt{q}.
\]
If \(\wt{D}\in I_p\), uniqueness of Euclidean division by \(p\) yields the final
statement.
\end{proof}

\begin{corollary}[Exclusion of the \(p^{-1}\)-witness class under quotient range binding]
\label{cor:eliminate-pinv}
If a quotient cell is constrained by \eqref{eq:q31-recompose}--\eqref{eq:q31-lookups} or by \eqref{eq:q66-recompose}--\eqref{eq:q66-lookups}, then it cannot equal \(-p^{-1}\pmod r\). Hence the specific \(p^{-1}\)-based vacuity attack is excluded.
\end{corollary}

\begin{proof}
Any such quotient cell is an integer \(\wt q\) in \([0,2^{31})\) or in
\([0,2^{66})\). Assume for contradiction that
\[
\wt q \equiv -p^{-1}\pmod r.
\]
Then
\[
p\wt q + 1 \equiv 0 \pmod r.
\]
If \(\wt q<2^{31}\), then
\[
0 < p\wt q+1 \le p(2^{31}-1)+1 < p\cdot 2^{31}+p < r.
\]
If \(\wt q<2^{66}\), then
\[
0 < p\wt q+1 \le p(2^{66}-1)+1 < p\cdot 2^{66}+p < r.
\]
In either case \(0<p\wt q+1<r\), so it cannot be divisible by \(r\), a
contradiction. Therefore \(\wt q\not\equiv -p^{-1}\pmod r\).
\end{proof}

\subsection{Exact row soundness}

For each active row relation, once a local realization is matched to an audited
normalized template and all participating cells are canonically range-bound, we
write
\[
\wt{L_t[j]},\qquad \wt{D_t[j]},\qquad \wt{Q_t[j]}
\]
for the corresponding audited/canonical integer evaluations. Concretely, on
active rows these are the integers supplied by
Definition~\ref{def:template-realization-coherence} together with
Lemma~\ref{lem:exact-audited-template-evaluation}; on inactive rows they are
zero by Definition~\ref{def:rowtable-zero-extension}. Their images in \(R\) are
denoted
\[
L_t[j],\qquad D_t[j],\qquad Q_t[j].
\]
We enforce local relations as
\[
L_t[j]=D_t[j]+pQ_t[j]\quad\text{in }R.
\]

For example:
\begin{itemize}[leftmargin=2em]
    \item inverse row:
    \[
    L = q\cdot d;
    \]
    \item \(E\times F\) coordinate row:
    \[
    L = \wt{x}_k\,\wt{a};
    \]
    \item \(E\times E\) coordinate row:
    \[
    L = \sum_{i,j=0}^{3} m_{ij}^k\,\wt{x}_i\,\wt{y}_j.
    \]
\end{itemize}
Here \(m_{ij}^k\) denotes an audited non-negative normalized template
coefficient in the instantiated library, not a raw signed basis coefficient.

To upgrade field equalities to integer equalities in a soundness proof, we require
the range bound for \emph{all satisfying witnesses} (not only honest ones).

\begin{definition}[Row-range discipline]
\label{def:row-range-discipline}
For each local relation \(t\), there exists \(B_t\in\{2^{31},2^{66}\}\) such that
every satisfying row assignment obeys
\[
0 \le \wt{L_t[j]} < pB_t + p < r.
\]
\end{definition}

\begin{lemma}[Monomial no-wrap]
\label{lem:monomial-nowrap}
Let \(x,y\in R\) satisfy:
\begin{enumerate}[label=(\roman*),leftmargin=2em]
    \item \(x\) is the image in \(R\) of an integer \(\wt{x}\) with
    \(0\le \wt{x}\le U_x\);
    \item \(y\) is the image in \(R\) of an integer \(\wt{y}\) with
    \(0\le \wt{y}\le U_y\);
    \item \(U_xU_y<r\).
\end{enumerate}
Then the field product \(xy\in R\) is the image in \(R\) of the integer product
\(\wt{x}\wt{y}\), and \(0\le \wt{x}\wt{y}<r\).
\end{lemma}

\begin{proof}
In \(\Z/r\Z\) we have \(x\equiv \wt{x}\) and \(y\equiv \wt{y}\), hence
\(xy\equiv \wt{x}\wt{y}\pmod r\). By (iii),
\(0\le \wt{x}\wt{y}\le U_xU_y<r\), so reduction modulo \(r\) is trivial and
\(xy\) is represented by \(\wt{x}\wt{y}\) in \(R\).
\end{proof}

\paragraph{Audited no-wrap bounds (moved to appendix).}
We discharge Definition~\ref{def:row-range-discipline} and the corresponding
RHS no-wrap bounds for all satisfying witnesses using a public template audit
certificate and a deterministic bound compiler; see
Appendix~\ref{app:audit-certificate}.

\paragraph{Template convention (non-negative normalization).}
This note (and the public audit certificate in Appendix~\ref{app:audit-certificate})
restricts to row templates already compiled into a \emph{non-negative} deferred-quotient
form \(L_t = D_t + pQ_t\), where \(L_t\) and \(D_t\) are non-negative combinations
of degree-\(\le 2\) monomials over ranged row variables and \(Q_t\) is bound to one
of the declared non-negative quotient intervals (e.g., \(I_{31}\) or \(I_{66}\)).
Handling row relations whose satisfying assignments force a negative Euclidean
quotient (e.g., \(L_t < D_t\) on some satisfying rows) would require signed quotient
ranges or a different normalization scheme, and is outside the scope of this
repair note.

\paragraph{Worked normalization example.}
A simple illustration is a row relation initially written as
\[
x_0+x_1-x_2 \equiv 0 \pmod p,
\]
where \(x_0,x_1,x_2\) are designated residue cells. In the restricted model of
this note, such a relation is compiled to the non-negative deferred-quotient
form
\[
x_0+x_1 = x_2 + p q,
\]
with \(x_2\in I_p\) and \(q\in\{0,1\}\), because
\[
0\le x_0+x_1\le 2(p-1)<2p.
\]
This is the kind of normalization covered by the present repair note. By
contrast, families whose satisfying assignments force negative quotients or a
signed remainder convention require a different normalization scheme and are
outside the present scope.

\paragraph{Horner-parameter convention.}
Throughout this note, the promoted-Horner parameter \(\rho\) is a fixed public
extension-field constant \(\rho\in\Ep\). It is not a witness cell and not an
instance-dependent Fiat--Shamir challenge. The published
audit/certificate payload is tied to this pinned value; changing \(\rho\)
requires regenerating the promoted-Horner template specialization and its bound
certificate payload.

\begin{theorem}[Core rowwise exactness from canonical lifts]
\label{thm:core-rowwise-exactness}
Fix a satisfied row relation
\[
L_t[j]=D_t[j]+pQ_t[j]\quad\text{in }R.
\]
Assume there exists \(B_t\in\{2^{31},2^{66}\}\) and integers
\(\wt{L_t[j]},\wt{D_t[j]},\wt{Q_t[j]}\) such that:
\begin{enumerate}[label=(\roman*),leftmargin=2em]
    \item \(L_t[j]\), \(D_t[j]\), and \(Q_t[j]\) are the images in \(R\) of
    \(\wt{L_t[j]}\), \(\wt{D_t[j]}\), and \(\wt{Q_t[j]}\), respectively;
    \item
    \[
    0\le \wt{L_t[j]}<r;
    \]
    \item
    \[
    0\le \wt{Q_t[j]}<B_t;
    \]
    \item
    \[
    0\le \wt{D_t[j]}+p\,\wt{Q_t[j]}<r.
    \]
\end{enumerate}
Then
\[
\wt{L_t[j]}=\wt{D_t[j]}+p\,\wt{Q_t[j]}
\quad\text{in }\Z.
\]
If additionally \(\wt{D_t[j]}\in I_p\), then
\[
\wt{D_t[j]}=\wt{L_t[j]}\bmod p,
\qquad
\wt{Q_t[j]}=\left\lfloor\frac{\wt{L_t[j]}}{p}\right\rfloor.
\]
\end{theorem}

\begin{proof}
By assumption (i), the row values are the images in \(R\) of the integers
\(\wt{L_t[j]},\wt{D_t[j]},\wt{Q_t[j]}\). Since
\[
L_t[j]=D_t[j]+pQ_t[j]\quad\text{in }R,
\]
we obtain
\[
\wt{L_t[j]} \equiv \wt{D_t[j]} + p\,\wt{Q_t[j]} \pmod r.
\]
By assumptions (ii) and (iv), both integers
\[
\wt{L_t[j]}
\qquad\text{and}\qquad
\wt{D_t[j]} + p\,\wt{Q_t[j]}
\]
lie in the interval \([0,r)\). Therefore they are equal in \(\Z\):
\[
\wt{L_t[j]}=\wt{D_t[j]}+p\,\wt{Q_t[j]}.
\]
If additionally \(\wt{D_t[j]}\in I_p\), uniqueness of Euclidean division by
\(p\) yields
\[
\wt{D_t[j]}=\wt{L_t[j]}\bmod p,
\qquad
\wt{Q_t[j]}=\left\lfloor\frac{\wt{L_t[j]}}{p}\right\rfloor.
\]
\end{proof}

\begin{theorem}[Certified rowwise no-vacuity in the restricted model]
\label{thm:rowwise-no-vacuity}
Assume the following four premises:
\begin{enumerate}[label=(\Alph*),leftmargin=2em]
    \item \textbf{Canonical binding.} Residue, quotient, and carry columns are
    range-bound by the binding constraints
    \eqref{eq:residue-recompose}--\eqref{eq:residue-not-p},
    \eqref{eq:q31-recompose}--\eqref{eq:q31-lookups},
    \eqref{eq:q66-recompose}--\eqref{eq:q66-lookups}, and
    \eqref{eq:bit-carry} or \eqref{eq:u8-lookup}. Canonical residue rows are
    exactly those rows whose \(D_t[j]\) is a designated residue witness cell
    \(c_t[j]\in I_p\).
    \item \textbf{Audited no-wrap bounds.} For every active matched family \(f\),
    the published audit certificate yields
    \[
    0\le \wt{L_t[j]}\le U_f^{\mathrm{audit}} < pB_f+p<r
    \]
    and
    \[
    0\le \wt{D_t[j]}\le U_f^{D,\mathrm{audit}},
    \qquad
    U_f^{D,\mathrm{audit}}+p(B_f-1)<r.
    \]
    In the released artifact/checker flow, the numerical and metadata side of
    this premise is mechanically verified from the public certificate.
    \item \textbf{Declared template/realization coherence.} On every active row,
    the local realization expressions agree with the assigned audited
    normalized template as specified by
    Definition~\ref{def:template-realization-coherence}. This coherence bridge
    is a published symbolic proof obligation; it is not fully mechanized by the
    released checker. A worked end-to-end symbolic coherence derivation is given
    for an \(E\times E\) row in Appendix~\ref{app:coverage-and-wiring}; the
    remaining instantiated families (\(\mathrm{inv},\mathrm{EF},\mathrm{hor},\mathrm{car}\))
    are structurally simpler and covered by the same symbolic certificate relation.
    \item \textbf{Complete wiring/coverage.} Every instantiated
    quotient/residue/carry slot is covered exactly once by the intended binding
    rules (no omissions), as formalized by
    Definitions~\ref{def:complete-quotient-coverage},
    \ref{def:selector-quotient-wiring}, and
    \ref{def:selector-residue-carry-wiring}. In the released artifact/checker
    flow, this coverage/wiring premise is mechanically verified from the public
    manifest/certificate metadata.
\end{enumerate}
Then for every satisfied row relation
\begin{equation}
L_t[j] = D_t[j] + pQ_t[j]
\quad\text{in }R,
\label{eq:row-relation-R}
\end{equation}
the hypotheses of Theorem~\ref{thm:core-rowwise-exactness} hold. Consequently,
\begin{equation}
\wt{L_t[j]} = \wt{D_t[j]} + p\,\wt{Q_t[j]}
\quad\text{in }\Z.
\label{eq:row-relation-Z}
\end{equation}
For canonical residue rows, \(c_t[j]\) is the canonical remainder modulo \(p\),
and \(Q_t[j]\) is the exact Euclidean quotient.
\end{theorem}

\begin{proof}
Fix any row relation \((t,j)\) satisfying \eqref{eq:row-relation-R}. By
canonical binding together with
Lemmas~\ref{lem:no-omitted-quotient-wire} and
\ref{lem:no-omitted-residue-carry-wire}, every quotient/residue/carry cell
participating in that row admits a canonical integer lift in its declared
interval.

If \(j\notin\Omega_t\), then Definition~\ref{def:rowtable-zero-extension} gives
\[
L_t[j]=D_t[j]=Q_t[j]=0,
\]
so the hypotheses of Theorem~\ref{thm:core-rowwise-exactness} are immediate.

Assume now \(j\in\Omega_t\). By
Definition~\ref{def:template-realization-coherence} and
Lemma~\ref{lem:exact-audited-template-evaluation}, the values \(L_t[j]\) and
\(D_t[j]\) are the images in \(R\) of exact audited integer evaluations. By
Lemma~\ref{lem:row-range-from-audit},
\[
0\le \wt{L_t[j]} < pB_t+p<r
\]
for the class \(B_t\) of the matched family. By
Lemma~\ref{lem:no-omitted-quotient-wire} together with
Definition~\ref{def:family-quotient-class-coherence},
\[
0\le \wt{Q_t[j]}<B_t.
\]
By Lemma~\ref{lem:rhs-nowrap-from-audit},
\[
0\le \wt{D_t[j]}+p\wt{Q_t[j]}<r.
\]
Thus all hypotheses of Theorem~\ref{thm:core-rowwise-exactness} hold, yielding
\eqref{eq:row-relation-Z}. On canonical residue rows, the designated residue
cell \(c_t[j]=D_t[j]\) lies in \(I_p\), so the final Euclidean-division
statement follows from Theorem~\ref{thm:core-rowwise-exactness}.
\end{proof}

\paragraph{Reading of the main theorem.}
Theorem~\ref{thm:rowwise-no-vacuity} should therefore be read as a certified
restricted-model theorem: under the stated binding, audit, wiring, and
coherence premises, it upgrades satisfied deferred-quotient row identities to
exact \(\Z\)-equalities and excludes the specific rowwise
\(p^{-1}\)-witness vacuity class in this model.

\subsection{Extension to a concrete Halo2-style backend (discussion)}

The Tier-2 wrapper check itself can remain unchanged; typically it enforces the
single compressed evaluation identity
\begin{equation}
\sum_t \lambda_t\bigl(l_t-d_t-pq_t\bigr) = h\,Z_\Omega(\zeta),
\label{eq:main-compressed-check}
\end{equation}
where
\[
l_t = L_t(\zeta),\quad
d_t = D_t(\zeta),\quad
q_t = Q_t(\zeta),\quad
h = H(\zeta),
\]
and \(\zeta \in R\setminus \Omega\) is a verifier challenge (sampled interactively
or via Fiat--Shamir).
The identification between compression-layer values
\(L_t(\omega^j),D_t(\omega^j)\) and rowwise values \(L_t[j],D_t[j]\) is fixed
by Definition~\ref{def:compression-rowwise-rhs-consistency}.

\begin{remark}[Backend-dependent closure boundary]
\label{rem:wrapper-extension}
Theorem~\ref{thm:rowwise-no-vacuity} is the core algebraic repair in this note.
A fully closed cryptographic wrapper theorem additionally requires a separate
backend proof of transcript binding/extraction, degree discipline for
\eqref{eq:main-compressed-check}, and Fiat--Shamir soundness.
This note does not prove such a backend-closure theorem.
Instead, the companion artifact provides a pinned backend manifest only for
reproducible certificate/checker binding, while full backend closure remains a
separate proof obligation.
\end{remark}

\subsection{Tier-2 circuit cost accounting}

The main arithmetic compression check \eqref{eq:main-compressed-check} is unchanged. The only additional cost is the canonical-binding layer above.
Unless stated otherwise, all symbolic counts in this subsection are exact arithmetic-object counts under this
bookkeeping model (lookup cells, native multiplicative constraints, and linear
recomposition constraints). They do not claim exact wall-clock proving time,
which depends on backend implementation details.
In particular, this is not a claim that one lookup cell is cost-equivalent to
one native multiplication gate; we report them as distinct arithmetic objects.
\paragraph{Supplementary implementation note (not a main claim).}
An implemented Halo2 B1-lite circuit-shape comparison is reported only as a
supplementary reproducibility artifact in
Appendix~\ref{app:b1lite-supplement}. It is intentionally not
task-equated (baseline omits modular reduction and CRT-consistency) and
is not used as a main quantitative claim.

Let
\begin{align*}
N_p &:= \#\{\text{residue cells bound by } \eqref{eq:residue-recompose}\text{--}\eqref{eq:residue-not-p}\},\\
N_{31} &:= \#\{\text{31-bit quotient cells bound by } \eqref{eq:q31-recompose}\text{--}\eqref{eq:q31-lookups}\},\\
N_{66} &:= \#\{\text{66-bit quotient cells bound by } \eqref{eq:q66-recompose}\text{--}\eqref{eq:q66-lookups}\},\\
N_{\mathrm{bit}} &:= \#\{\text{boolean carry cells}\},\\
N_8 &:= \#\{\text{3-bit carry cells}\}.
\end{align*}
We decompose quotient-cell totals by family as
\[
N_{31}=N_{31}^{\mathrm{EF}}+N_{31}^{\mathrm{car}},
\qquad
N_{66}=N_{66}^{\mathrm{inv}}+N_{66}^{\mathrm{EE}}+N_{66}^{\mathrm{hor}},
\]
where \(N_{31}^{\mathrm{car}}\) is exactly the contribution of
\(Q_{\mathrm{car}}\) carry-normalization rows.

Then the added Tier-2 arithmetization cost is:

\paragraph{Lookup cells.}
\begin{equation}
L_{\mathrm{canon}}
=
2N_p + 2N_{31} + 5N_{66} + N_8.
\label{eq:canon-lookup-cost}
\end{equation}

\paragraph{Native multiplicative constraints.}
\begin{equation}
M_{\mathrm{canon}}
=
N_p + N_{\mathrm{bit}}.
\label{eq:canon-mul-cost}
\end{equation}
These are:
\begin{itemize}[leftmargin=2em]
    \item \(N_p\) non-equality constraints \((x-p)\nu_x=1\),
    \item \(N_{\mathrm{bit}}\) boolean constraints \(c(c-1)=0\).
\end{itemize}

\paragraph{Linear recomposition constraints.}
\begin{equation}
\mathrm{Lin}_{\mathrm{canon}}
=
N_p + N_{31} + N_{66}.
\label{eq:canon-lin-cost}
\end{equation}

No limb-wise multiplication, no CRT consistency, and no foreign-field carry propagation appear anywhere in the canonical-binding layer.

\subsubsection{Concrete promoted Horner row (scope-separated from carry-normalization rows)}

Consider one promoted Horner row:
\[
s = 2a_0+a_1+a_2+a_3,
\qquad
s' = a_0 + \rho(a_1+\rho(a_2+\rho a_3)),
\]
where \(\rho\) is the fixed Horner scalar parameter of the instantiated row
family, with \(\rho\in\Ep\) pinned for the instantiated promoted-Horner family
(not a witness variable and not an instance-dependent Fiat--Shamir challenge).
Multiplications by \(\rho\) are compiled into fixed specialized templates
recorded in the public audit payload.
With exact intermediate variables
\[
m_3,\; h_2,\; m_2,\; h_1,\; m_1
\]
and quotient/carry vectors as in the deferred-quotient arithmetization.
This micro-count intentionally excludes separate carry-normalization rows.
Although the underlying arithmetic corresponds to one Horner evaluation step of
a degree-\(3\) expression, the circuit realizes it via a fixed family of
linear/quadratic subtemplates with intermediate witnesses; hence the audited
template language remains degree-\(\le 2\).
All Horner-related audit bounds and quotient-class assignments refer to these
subtemplates, not to a single degree-\(3\) monomial template.
Then one promoted Horner row contains:
\begin{itemize}[leftmargin=2em]
    \item \(11\) \(E\)-elements of canonical residues,
    \[
    a_0,a_1,a_2,a_3,s,s',m_3,h_2,m_2,h_1,m_1,
    \]
    hence
    \[
    N_p^{\mathrm{round}} = 11\cdot 4 = 44;
    \]
    \item \(3\) quotient vectors \(q_3,q_2,q_1\), hence
    \[
    N_{66}^{\mathrm{round}} = 3\cdot 4 = 12;
    \]
    \item carry vectors \(c_2,c_1,v\in\{0,1\}^4\), hence
    \[
    N_{\mathrm{bit}}^{\mathrm{round}} = 12;
    \]
    \item one 3-bit carry vector \(u\in\{0,\dots,4\}^4\), hence
    \[
    N_8^{\mathrm{round}} = 4.
    \]
\end{itemize}
In this scope,
\[
N_{31}^{\mathrm{round}}=0,
\]
because no \(Q_{\mathrm{car}}\) row is included.

Therefore the added binding counts per promoted Horner row are exactly:
\begin{equation}
L_{\mathrm{canon}}^{\mathrm{round}}
=
2\cdot 44 + 5\cdot 12 + 4
=
152
\label{eq:round-lookup-cost}
\end{equation}
lookup cells,
\begin{equation}
M_{\mathrm{canon}}^{\mathrm{round}}
=
44 + 12
=
56
\label{eq:round-mul-cost}
\end{equation}
native multiplicative constraints, and
\begin{equation}
\mathrm{Lin}_{\mathrm{canon}}^{\mathrm{round}}
=
44 + 12
=
56
\label{eq:round-lin-cost}
\end{equation}
linear recomposition constraints.

\paragraph{Carry-normalization add-on (31-bit quotient family).}
A carry-normalization row contributes one quotient vector
\(\{q_{\mathrm{car}}^{(k)}\}_{k=0}^3\), i.e.
\[
N_{31,\mathrm{car}}^{\mathrm{row}}=4.
\]
Hence each such row adds
\[
\Delta L_{\mathrm{canon}}=2\cdot 4=8,\qquad
\Delta M_{\mathrm{canon}}=0,\qquad
\Delta \mathrm{Lin}_{\mathrm{canon}}=4.
\]
If one carry-normalization row is paired with each promoted Horner row, the
per-package counts are
\[
L_{\mathrm{canon}}^{\mathrm{pkg}}=160,\qquad
M_{\mathrm{canon}}^{\mathrm{pkg}}=56,\qquad
\mathrm{Lin}_{\mathrm{canon}}^{\mathrm{pkg}}=60.
\]
The carry-normalization add-on reuses the carry witnesses already counted in
the paired promoted Horner row and introduces only the fresh quotient vector
\(\{q_{\mathrm{car}}^{(k)}\}_{k=0}^3\) as additional bounded cells. Hence it
contributes no extra boolean or \(3\)-bit carry cost in this package
accounting.

Let \(R_{\mathrm{hor}}\) be the number of promoted Horner rows and
\(R_{\mathrm{car}}\) the number of carry-normalization rows. Then the total added
binding overhead from these two row families is
\begin{equation}
L_{\mathrm{canon}}^{\mathrm{tot}}
=
152\,R_{\mathrm{hor}}+8\,R_{\mathrm{car}},
\label{eq:total-lookup-cost}
\end{equation}
\begin{equation}
M_{\mathrm{canon}}^{\mathrm{tot}}
=
56\,R_{\mathrm{hor}},
\label{eq:total-mul-cost}
\end{equation}
and
\begin{equation}
\mathrm{Lin}_{\mathrm{canon}}^{\mathrm{tot}}
=
56\,R_{\mathrm{hor}}+4\,R_{\mathrm{car}}.
\label{eq:total-lin-cost}
\end{equation}
These three equations are exact symbolic arithmetic-object counts for the added
canonical-binding layer in this note.
This is consistent with
\eqref{eq:canon-lookup-cost}--\eqref{eq:canon-lin-cost} under
\[
N_{66}^{\mathrm{hor}}=12R_{\mathrm{hor}},
\qquad
N_{31}^{\mathrm{car}}=4R_{\mathrm{car}}.
\]

\paragraph{Instantiation for standard \(n\)-variable transcripts.}
For the common fractional-GKR transcript family, if carry-normalization rows
are instantiated one-for-one with promoted Horner rows
\[
R_{\mathrm{car}}=R_{\mathrm{hor}}=:R_{\mathrm{sc}},
\]
and
\[
R_{\mathrm{sc}}
=
\sum_{k=1}^{n-1} k
=
\frac{n(n-1)}{2},
\]
equations \eqref{eq:total-lookup-cost}--\eqref{eq:total-lin-cost} become
\[
L_{\mathrm{canon}}^{\mathrm{tot}}=80\,n(n-1),
\qquad
M_{\mathrm{canon}}^{\mathrm{tot}}=28\,n(n-1),
\qquad
\mathrm{Lin}_{\mathrm{canon}}^{\mathrm{tot}}=30\,n(n-1).
\]

\paragraph{Asymptotic position versus full emulation (contextual corollary).}
Let \(N_{\mathrm{base}}\) denote the size of the original computation being
wrapped, and \(m:=|\Omega|\) the wrapped transcript-table size. In the standard
fractional-GKR schedule targeted by the companion construction, we adopt
\[
m=\Theta(\log^2 N_{\mathrm{base}}),
\qquad
R_{\mathrm{sc}}=\Theta(m).
\]
This note does not re-prove these schedule scalings; it uses them as contextual
input to interpret the exact symbolic object counts above. Under this adopted
schedule, the canonical-binding overhead is
\[
\Theta(R_{\mathrm{sc}})
=
\Theta(m)
=
\Theta(\log^2 N_{\mathrm{base}}).
\]
By contrast, full in-circuit replay of the original computation costs
\(\Omega(N_{\mathrm{base}})\). Therefore the overhead ratio is
\[
O\!\left(\frac{\log^2 N_{\mathrm{base}}}{N_{\mathrm{base}}}\right),
\]
which vanishes as \(N_{\mathrm{base}}\) grows.

\subsubsection{Representative symbolic instantiations and cell-level baseline}

The formulas above already give an exact accounting of the \emph{added} binding
overhead in terms of arithmetic objects. This equation-level accounting
(\eqref{eq:total-lookup-cost}--\eqref{eq:total-lin-cost}) is the primary
quantitative claim of this repair note.

\begin{table}[t]
\centering
\small
\setlength{\tabcolsep}{6pt}
\begin{tabular}{|r|r|r|r|r|}
\hline
\(n\) & \(R_{\mathrm{sc}}=\frac{n(n-1)}{2}\) & \(L_{\mathrm{canon}}^{\mathrm{tot}}\) & \(M_{\mathrm{canon}}^{\mathrm{tot}}\) & \(\mathrm{Lin}_{\mathrm{canon}}^{\mathrm{tot}}\) \\
\hline
16  & 120  & 19,200    & 6,720    & 7,200 \\
32  & 496  & 79,360    & 27,776   & 29,760 \\
64  & 2,016 & 322,560   & 112,896  & 120,960 \\
128 & 8,128 & 1,300,480 & 455,168  & 487,680 \\
\hline
\end{tabular}
\caption{Added canonical-binding overhead for representative \(n\) under the
standard schedule \(R_{\mathrm{car}}=R_{\mathrm{hor}}=R_{\mathrm{sc}}\).}
\label{tab:canon-overhead-bench}
\end{table}

\paragraph{Binding primitive comparison (lookup limbs vs bit decomposition).}
As a backend-agnostic qualitative baseline at the cell level, we compare our
lookup-limb binding against explicit bit decomposition.
As a simple baseline, one can range-bind a \(k\)-bit value by explicit bit
decomposition (introduce \(k\) boolean variables plus one recomposition
constraint). Our binding mechanism replaces \(\Theta(k)\) multiplicative
constraints per cell with \(O(1)\) lookup limbs and one recomposition, under
the explicit premise that the Tier-2 backend provides a native lookup argument.

\begin{table}[t]
\centering
\small
\setlength{\tabcolsep}{6pt}
\begin{tabular}{|l|l|l|}
\hline
Cell type & This note (lookup-limb binding) & Bit-decomposition baseline \\
\hline
Residue cell in \(I_p\) & \(2\) lookups + \(1\) linear + \(1\) non-equality & \(31\) booleans + \(1\) linear + \(1\) non-equality \\
31-bit quotient cell in \(I_{31}\) & \(2\) lookups + \(1\) linear & \(31\) booleans + \(1\) linear \\
66-bit quotient cell in \(I_{66}\) & \(5\) lookups + \(1\) linear & \(66\) booleans + \(1\) linear \\
Boolean carry cell & \(1\) boolean & \(1\) boolean \\
3-bit carry cell in \(T_8\) & \(1\) lookup & \(3\) booleans + \(1\) linear \\
\hline
\end{tabular}
\caption{Binding cost per cell (under native lookup availability): lookup limbs
versus explicit bit decomposition.}
\label{tab:binding-vs-bitdecomp}
\end{table}

For example, in the promoted Horner scope above with
\(N_p^{\mathrm{round}}=44\),
\(N_{66}^{\mathrm{round}}=12\),
\(N_{\mathrm{bit}}^{\mathrm{round}}=12\), and
\(N_8^{\mathrm{round}}=4\),
explicit bit decomposition would add
\[
31N_p^{\mathrm{round}} + 66N_{66}^{\mathrm{round}} + N_{\mathrm{bit}}^{\mathrm{round}} + 3N_8^{\mathrm{round}} + N_p^{\mathrm{round}} = 2224
\]
native multiplicative constraints (booleanity plus residue non-equality),
whereas the lookup-limb package adds
\(M_{\mathrm{canon}}^{\mathrm{round}}=56\) multiplicative constraints and
\(L_{\mathrm{canon}}^{\mathrm{round}}=152\) lookup cells.

\paragraph{Implementation note (supplementary B1-lite artifact only).}
For reproducibility, we report one implemented Halo2 circuit-shape comparison
in Appendix~\ref{app:b1lite-supplement}. This supplementary artifact is
backend-specific, not task-equated, and not part of the main theorem
claims or main quantitative evidence.

\subsection{Summary}

The simple algebraic repair in this model to eliminate the \(p^{-1}\)-based vacuity is:
\begin{enumerate}[label=(\arabic*),leftmargin=2em]
    \item retain the one-limb exact arithmetic over \(R\),
    \item keep the single compressed polynomial evaluation check \eqref{eq:main-compressed-check},
    \item add only a \emph{native} canonical-binding layer consisting of:
    \begin{itemize}[leftmargin=2em]
        \item \(31\)-bit lookup \(+\) one non-equality-to-\(p\) gate for each residue cell,
        \item \(31\)-bit or \(66\)-bit power-of-two lookup bounds for quotient cells,
        \item boolean or \(3\)-bit lookup constraints for carries.
    \end{itemize}
\end{enumerate}

Within the restricted certified model above, this excludes the
\(p^{-1}\)-based vacuity class and realizes the missing canonical-lift property
guarantee \(\mathcal{P}_{\mathrm{canon}}\) at the row level while keeping Tier-2
arithmetic fully native over \(\F_r\). The additional wrapper
arithmetic-object count in this bookkeeping model is still given by
\eqref{eq:total-lookup-cost}, \eqref{eq:total-mul-cost}, and
\eqref{eq:total-lin-cost}, with representative instantiations in
Table~\ref{tab:canon-overhead-bench} and a cell-level qualitative comparison in
Table~\ref{tab:binding-vs-bitdecomp}.

At claim level, the note proves a certified rowwise exactness theorem for this
restricted model. The released checker machine-verifies
bounds/wiring/metadata/digest integrity, while full symbolic
template/realization coherence remains a published manuscript-level proof
obligation. Accordingly, this note does not by itself establish a
backend-independent wrapper soundness theorem; wrapper-level cryptographic
closure remains the separate backend-dependent task described in
Remark~\ref{rem:wrapper-extension}.

\appendix

% BEGIN OPTION2 APPENDICES
\section{Audited no-wrap bounds and template certificates}
\label{app:audit-certificate}
This appendix records the public audit certificate, deterministic bound
verification, and instantiated constants used to justify the no-wrap
premises in the rowwise exactness layer.

\paragraph{Row-family bound audit (all satisfying witnesses).}
To avoid treating Definition~\ref{def:row-range-discipline} as a standalone
assumption, we publish the following row-family audit. Each row template is
encoded in normalized non-negative deferred-quotient form and assigned a
class \(B_f\in\{2^{31},2^{66}\}\).

\begin{center}
\begin{tabular}{|l|l|l|l|}
\hline
Row family \(f\) & Integer form for \(L_f\) & Certified bound \(U_f^{\mathrm{audit}}\) & Class \(B_f\) \\
\hline
inverse & \(L=q\cdot d\) & \((2^{66}-1)(p-1)\) & \(2^{66}\) \\
\hline
\(E\times F\) coordinate & \(L=\wt{x}_k\,\wt{a}\) & \((p-1)^2\) & \(2^{31}\) \\
\hline
\(E\times E\) coordinate & \(L=\sum_{i,j} m_{ij}^k \wt{x}_i\wt{y}_j\) & \(16\,(p-1)^3\) & \(2^{66}\) \\
\hline
promoted Horner & family of linear/quadratic subtemplates (promoted Horner block) & \(U_{\mathrm{hor}}^{\mathrm{audit}}\) & \(2^{66}\) \\
\hline
carry normalization & linear carry-transfer form & \(U_{\mathrm{car}}^{\mathrm{audit}}\) & \(2^{31}\) \\
\hline
\end{tabular}
\end{center}

Here
\[
M_{\mathrm{EE}}^{\max}
:=
\max_{k,i,j}\left|\widetilde{m_{ij}^k}\right|
\le
p-1
=
2147483646
\]
is a conservative worst-case bound under canonical integer lifts (independent
of basis coefficients).

\begin{definition}[Template-family audit certificate and deterministic bound compiler]
\label{def:template-audit-certificate}
A public row-audit certificate consists of a finite family index set
\(\mathcal F\). Each family \(f\in\mathcal F\) has a finite template set
\(\mathcal T_f\). For each template \(\tau\in\mathcal T_f\), the certificate
provides:
\begin{enumerate}[label=(\roman*),leftmargin=2em]
    \item a normalized non-negative left polynomial
    \[
    L_{f,\tau}(\mathbf z)
    =
    \sum_{u\in\mathcal L_{f,\tau}} a_{f,\tau,u}\,z_u
    +
    \sum_{(u,v)\in\mathcal B_{f,\tau}} b_{f,\tau,uv}\,z_u z_v,
    \]
    with coefficients \(a_{f,\tau,u},b_{f,\tau,uv}\in\Z_{\ge 0}\);
    \item a normalized non-negative right polynomial
    \[
    D_{f,\tau}(\mathbf z)
    =
    \sum_{u\in\mathcal L^D_{f,\tau}} \alpha_{f,\tau,u}\,z_u
    +
    \sum_{(u,v)\in\mathcal B^D_{f,\tau}} \beta_{f,\tau,uv}\,z_u z_v,
    \]
    with coefficients \(\alpha_{f,\tau,u},\beta_{f,\tau,uv}\in\Z_{\ge 0}\);
    \item interval bounds
    \[
    0 \le z_u \le U_{f,u}^{\max}
    \]
    induced by the class/range constraints;
    \item the class \(B_f\in\{2^{31},2^{66}\}\);
    \item deterministic per-template bounds
    \[
    U_{f,\tau}^{L}
    :=
    \sum_{u\in\mathcal L_{f,\tau}} a_{f,\tau,u}\,U_{f,u}^{\max}
    +
    \sum_{(u,v)\in\mathcal B_{f,\tau}} b_{f,\tau,uv}\,U_{f,u}^{\max}U_{f,v}^{\max},
    \]
    \[
    U_{f,\tau}^{D}
    :=
    \sum_{u\in\mathcal L^D_{f,\tau}} \alpha_{f,\tau,u}\,U_{f,u}^{\max}
    +
    \sum_{(u,v)\in\mathcal B^D_{f,\tau}} \beta_{f,\tau,uv}\,U_{f,u}^{\max}U_{f,v}^{\max}.
    \]
\end{enumerate}
Define family-level maxima
\[
U_f^{\mathrm{audit}}
:=
\max_{\tau\in\mathcal T_f} U_{f,\tau}^{L},
\qquad
U_f^{D,\mathrm{audit}}
:=
\max_{\tau\in\mathcal T_f} U_{f,\tau}^{D}.
\]
The compiler accepts only if each active instantiated row is syntactically
matched to exactly one template \((f,\tau)\) with \(\tau\in\mathcal T_f\),
while inactive row-table entries are handled separately by
Definition~\ref{def:rowtable-zero-extension}, and
\[
U_f^{\mathrm{audit}} < pB_f + p,
\qquad
U_f^{D,\mathrm{audit}} + p(B_f-1) < r
\]
holds for every \(f\). For the promoted-Horner and carry-normalization
families in the audit table above, we write
\[
U_{\mathrm{hor}}^{\mathrm{audit}}
:=
U_{f_{\mathrm{hor}}}^{\mathrm{audit}},
\qquad
U_{\mathrm{car}}^{\mathrm{audit}}
:=
U_{f_{\mathrm{car}}}^{\mathrm{audit}}.
\]
\[
U_{\mathrm{hor}}^{D,\mathrm{audit}}
:=
U_{f_{\mathrm{hor}}}^{D,\mathrm{audit}},
\qquad
U_{\mathrm{car}}^{D,\mathrm{audit}}
:=
U_{f_{\mathrm{car}}}^{D,\mathrm{audit}}.
\]
\end{definition}

\begin{definition}[Family/quotient-class coherence]
\label{def:family-quotient-class-coherence}
Whenever an active row is matched to template \((f,\tau)\) with family class
\(B_f\in\{2^{31},2^{66}\}\), every quotient slot appearing in that instantiated
template has local class tag
\[
b_{s,\ell}=
\begin{cases}
31,& B_f=2^{31},\\
66,& B_f=2^{66}.
\end{cases}
\]
Equivalently, every instantiated quotient value \(Q_t[j]\) in family \(f\) is
wired to a column in \(\mathcal Q_{31}\) if \(B_f=2^{31}\), and to a column in
\(\mathcal Q_{66}\) if \(B_f=2^{66}\).
\end{definition}

\paragraph{Symbolic certificate verifier relation (manuscript-level).}
Given the certificate data, the manuscript-level symbolic verifier relation
requires the following checks:
\begin{enumerate}[label=(\arabic*),leftmargin=2em]
    \item for each template \((f,\tau)\), parse \(L_{f,\tau}\) and \(D_{f,\tau}\),
    and ensure all certified coefficients are in \(\Z_{\ge 0}\);
    \item compute \(U_{f,\tau}^{L}\), \(U_{f,\tau}^{D}\), then family maxima
    \(U_f^{\mathrm{audit}}\), \(U_f^{D,\mathrm{audit}}\) from
    Definition~\ref{def:template-audit-certificate};
    \item check
    \[
    U_f^{\mathrm{audit}} < pB_f+p,
    \qquad
    U_f^{D,\mathrm{audit}} + p(B_f-1) < r
    \]
    for each family class \(B_f\);
    \item verify the template map \(\omega_T\) is total and functional on all
    active local scalar relation instances \((s,\rho,j)\), that each such
    instance is matched to exactly one template \((f,\tau)\) with
    \(\tau\in\mathcal T_f\), and that quotient-slot class tags are coherent
    with \(B_f\) as required by
    Definition~\ref{def:family-quotient-class-coherence}.
\end{enumerate}
These checks define the manuscript-level symbolic certificate relation and are
deterministic and independent of prover-provided runtime values.

\paragraph{Partially machine-checkable public certificate and checker.}
For this manuscript package we publish a concrete public
certificate
\texttt{certificates/public\_certificate.json}
and a deterministic checker
\texttt{scripts/check\_public\_certificate.py}, yielding a partially
machine-checkable certificate workflow.
The checker deterministically:
\begin{enumerate}[label=(\roman*),leftmargin=2em]
    \item reconstructs per-template bounds and family maxima;
    \item verifies the class/headroom checks
    \(U_f^{\mathrm{audit}}<pB_f+p\),
    \(U_f^{D,\mathrm{audit}}+p(B_f-1)<r\);
    \item verifies quotient/non-quotient coverage and selector-level wiring
    single-valued mapping/class consistency;
    \item verifies relation-template reference integrity
    (\texttt{relation\_id}\(\mapsto\)(\texttt{template\_family},
    \texttt{template\_id})), relation/row-activity metadata consistency, and
    inactive-row zero-extension flags;
    \item verifies digest binding between the certificate payload and the
    manuscript hash.
\end{enumerate}
In this release, the deterministic checker does \emph{not} mechanize full
symbolic equality validation of declared local realization expressions against
normalized template polynomials; that coherence relation remains specified as a
published symbolic certificate relation (Definition~\ref{def:template-realization-coherence}
and Lemma~\ref{lem:exact-audited-template-evaluation}).
Reproduction command:
\texttt{python scripts/check\_public\_certificate.py --certificate certificates/public\_certificate.json --manuscript wrapper\_note\_option2.tex}.
Pinned-instance check command:
\texttt{python scripts/check\_public\_certificate.py --certificate certificates/public\_certificate.json --manuscript wrapper\_note\_option2.tex --backend-instance certificates/h2dq\_backend\_instance.json}.
\paragraph{Backend-instance manifest role (reproducibility binding only).}
The optional backend-instance manifest binds the certificate payload digest,
checker hash, manuscript hash, and fixed backend schedule metadata for
reproducibility. In this note, it is used only for artifact pinning and
reproducible checker binding; it is not used as a standalone backend soundness
theorem.

\paragraph{Instantiated bounds for the current template library (current manuscript).}
\begin{center}
\begin{tabular}{|l|c|r|r|}
\hline
family id & class bits & decimal left bound & decimal right bound \\
\hline
\texttt{inv} & \(66\) & \(158456324880954722595264004098\) & \(2147483646\) \\
\hline
\texttt{EF} & \(31\) & \(4611686009837453316\) & \(2147483646\) \\
\hline
\texttt{EE} & \(66\) & \(158456324585806817830375522176\) & \(2147483646\) \\
\hline
\texttt{hor} & \(66\) & \(158456324585806817830375522176\) & \(158456324585806817830375522176\) \\
\hline
\texttt{car} & \(31\) & \(4294967299\) & \(4294967299\) \\
\hline
\end{tabular}
\end{center}
For this manuscript draft, the table above is encoded exactly in
\texttt{certificates/public\_certificate.json} (payload \texttt{families}) and
is checked against deterministic bound reconstruction by
\texttt{scripts/check\_public\_certificate.py}.
In this instantiation, inverse/\(E\times F\)/\(E\times E\) rows are canonical
residue rows, so \(D_t[j]=c_t[j]\in I_p\) and
\(U_{\mathrm{inv}}^{D,\mathrm{audit}}=
U_{\mathrm{EF}}^{D,\mathrm{audit}}=
U_{\mathrm{EE}}^{D,\mathrm{audit}}=p-1\).
In the current instantiated template library, the audited \(E\times E\) family
used for the constants above is treated in canonical non-negative form, so
\(D_t[j]=c_t[j]\in I_p\) on those rows. This note does not include alternative
signed normalizations that would require signed quotient handling.

\begin{lemma}[Template-derived family bounds]
\label{lem:template-bound-soundness}
Under Definition~\ref{def:template-audit-certificate}, for every family
\(f\in\mathcal F\), every template \(\tau\in\mathcal T_f\), and every satisfying
assignment \(\mathbf z\) of its row variables,
\[
0 \le L_{f,\tau}(\mathbf z) \le U_f^{\mathrm{audit}}.
\]
and
\[
0 \le D_{f,\tau}(\mathbf z) \le U_f^{D,\mathrm{audit}}.
\]
\end{lemma}

\begin{proof}
All certified coefficients are non-negative and every variable is restricted to
a non-negative interval. Therefore both \(L_{f,\tau}\) and \(D_{f,\tau}\) are
coordinate-wise monotone on the certified box
\(\prod_u [0,U_{f,u}^{\max}]\). Their maxima are attained at the upper endpoint
vector \((U_{f,u}^{\max})_u\), giving
\(U_f^{\mathrm{audit}}\) and \(U_f^{D,\mathrm{audit}}\) by definition.
Non-negativity gives the lower bounds.
\end{proof}

\paragraph{Concrete inequalities used in the audit.}
For \(p=2^{31}-1\):
\[
p\cdot 2^{66}+p
=
(2^{31}-1)(2^{66}+1)
=
2^{97}+2^{31}-2^{66}-1
<
2^{97}.
\]
The audit verifies:
\begin{align*}
&\text{inverse: } (2^{66}-1)(p-1)=158456324880954722595264004098 < p\cdot 2^{66}+p,\\
&E\times F: (p-1)^2=4611686009837453316 < p\cdot 2^{31}+p = 4611686018427387903,\\
&E\times E: 16\,M_{\mathrm{EE}}^{\max}(p-1)^2
=16(p-1)^3
=158456324585806817830375522176
< p\cdot 2^{66}+p.
\end{align*}
The compiler also checks
\[
U_{\mathrm{hor}}^{\mathrm{audit}} < p\cdot 2^{66}+p,
\qquad
U_{\mathrm{car}}^{\mathrm{audit}} < p\cdot 2^{31}+p.
\]
and the RHS-side no-wrap conditions
\[
U_{\mathrm{inv}}^{D,\mathrm{audit}}+p(2^{66}-1)<r,\quad
U_{\mathrm{EF}}^{D,\mathrm{audit}}+p(2^{31}-1)<r,\quad
U_{\mathrm{EE}}^{D,\mathrm{audit}}+p(2^{66}-1)<r,
\]
\[
U_{\mathrm{hor}}^{D,\mathrm{audit}}+p(2^{66}-1)<r,\quad
U_{\mathrm{car}}^{D,\mathrm{audit}}+p(2^{31}-1)<r.
\]

\paragraph{Appendix-style instantiated audit values for selected templates.}
For the current wrapper template set, the promoted-Horner and
carry-normalization family bounds are
instantiated as follows.
\begin{center}
\begin{tabular}{|l|l|l|}
\hline
Family & Instantiated template bound & Numeric value \\
\hline
promoted Horner &
\(
\max\!\{16M_{\mathrm{EE}}^{\max}(p-1)^2,\;2(p-1),\;5(p-1)\}
\) &
\(16(p-1)^3=158456324585806817830375522176\) \\
\hline
carry normalization &
\(
\max\!\{2(p-1)+7,\;2^{16}(2^{15}-1)+(2^{16}-1)\}
\) &
\(U_{\mathrm{car}}^{\mathrm{audit}}=4294967299\) \\
\hline
\end{tabular}
\end{center}

The promoted-Horner maximum comes from fixed-\(\rho\) multiplication sub-rows
\(m_3=\rho a_3\), \(m_2=\rho h_2\), \(m_1=\rho h_1\). In this note,
\(\rho\in\Ep\) is pinned at template-instantiation time, and these specialized
subtemplates are conservatively audited under the same bound regime as the
\(E\times E\) coordinate family. The linear Horner sub-rows
\(h_2=a_2+m_3\), \(h_1=a_1+m_2\), \(s'=a_0+m_1\), and
\(s=2a_0+a_1+a_2+a_3\) are strictly smaller.

The carry-normalization template family uses only low-width carry terms
(\(\{0,1\}\) and \(\{0,\dots,7\}\)) and 31-bit limb-recomposition terms, hence
the instantiated value above.

\begin{definition}[Instantiated full-family audit constants]
\label{def:instantiated-audit-constants}
For the current template set, define:
\begin{equation}
U_{\mathrm{inv}}^{\mathrm{audit}}
=
(2^{66}-1)(p-1)
=
158456324880954722595264004098,
\label{eq:u-inv-inst}
\end{equation}
\begin{equation}
U_{\mathrm{EF}}^{\mathrm{audit}}
=
(p-1)^2
=
4611686009837453316,
\label{eq:u-ef-inst}
\end{equation}
\begin{equation}
U_{\mathrm{EE}}^{\mathrm{audit}}
=
16(p-1)^3
=
158456324585806817830375522176,
\label{eq:u-ee-inst}
\end{equation}
\begin{equation}
U_{\mathrm{hor}}^{\mathrm{audit}}
=
16(p-1)^3
=
158456324585806817830375522176,
\label{eq:u-hor-inst}
\end{equation}
\begin{equation}
U_{\mathrm{car}}^{\mathrm{audit}}
=
4294967299.
\label{eq:u-car-inst}
\end{equation}
For RHS-side bounds, define:
\begin{equation}
U_{\mathrm{inv}}^{D,\mathrm{audit}}
=
p-1
=
2147483646,
\label{eq:u-inv-rhs-inst}
\end{equation}
\begin{equation}
U_{\mathrm{EF}}^{D,\mathrm{audit}}
=
p-1
=
2147483646,
\label{eq:u-ef-rhs-inst}
\end{equation}
\begin{equation}
U_{\mathrm{EE}}^{D,\mathrm{audit}}
=
p-1
=
2147483646,
\label{eq:u-ee-rhs-inst}
\end{equation}
\begin{equation}
U_{\mathrm{hor}}^{D,\mathrm{audit}}
=
16(p-1)^3
=
158456324585806817830375522176,
\label{eq:u-hor-rhs-inst}
\end{equation}
\begin{equation}
U_{\mathrm{car}}^{D,\mathrm{audit}}
=
4294967299.
\label{eq:u-car-rhs-inst}
\end{equation}
The family-to-bound/class map is fixed as
\[
(f_{\mathrm{inv}},B_{f_{\mathrm{inv}}},U_{f_{\mathrm{inv}}}^{\mathrm{audit}})
=(f_{\mathrm{inv}},2^{66},U_{\mathrm{inv}}^{\mathrm{audit}}),
\]
\[
(f_{\mathrm{EF}},B_{f_{\mathrm{EF}}},U_{f_{\mathrm{EF}}}^{\mathrm{audit}})
=(f_{\mathrm{EF}},2^{31},U_{\mathrm{EF}}^{\mathrm{audit}}),
\]
\[
(f_{\mathrm{EE}},B_{f_{\mathrm{EE}}},U_{f_{\mathrm{EE}}}^{\mathrm{audit}})
=(f_{\mathrm{EE}},2^{66},U_{\mathrm{EE}}^{\mathrm{audit}}),
\]
\[
(f_{\mathrm{hor}},B_{f_{\mathrm{hor}}},U_{f_{\mathrm{hor}}}^{\mathrm{audit}})
=(f_{\mathrm{hor}},2^{66},U_{\mathrm{hor}}^{\mathrm{audit}}),
\]
\[
(f_{\mathrm{car}},B_{f_{\mathrm{car}}},U_{f_{\mathrm{car}}}^{\mathrm{audit}})
=(f_{\mathrm{car}},2^{31},U_{\mathrm{car}}^{\mathrm{audit}}).
\]
\[
(f_{\mathrm{inv}},U_{f_{\mathrm{inv}}}^{D,\mathrm{audit}})
=(f_{\mathrm{inv}},U_{\mathrm{inv}}^{D,\mathrm{audit}}),
\]
\[
(f_{\mathrm{EF}},U_{f_{\mathrm{EF}}}^{D,\mathrm{audit}})
=(f_{\mathrm{EF}},U_{\mathrm{EF}}^{D,\mathrm{audit}}),
\]
\[
(f_{\mathrm{EE}},U_{f_{\mathrm{EE}}}^{D,\mathrm{audit}})
=(f_{\mathrm{EE}},U_{\mathrm{EE}}^{D,\mathrm{audit}}),
\]
\[
(f_{\mathrm{hor}},U_{f_{\mathrm{hor}}}^{D,\mathrm{audit}})
=(f_{\mathrm{hor}},U_{\mathrm{hor}}^{D,\mathrm{audit}}),
\]
\[
(f_{\mathrm{car}},U_{f_{\mathrm{car}}}^{D,\mathrm{audit}})
=(f_{\mathrm{car}},U_{\mathrm{car}}^{D,\mathrm{audit}}).
\]
\end{definition}

\begin{lemma}[Instantiated full-family bound closure]
\label{lem:instantiated-audit-closure}
Under Definition~\ref{def:instantiated-audit-constants},
\begin{equation}
U_{\mathrm{inv}}^{\mathrm{audit}} < p\cdot 2^{66}+p,
\label{eq:u-inv-inst-bound}
\end{equation}
\begin{equation}
U_{\mathrm{EF}}^{\mathrm{audit}} < p\cdot 2^{31}+p,
\label{eq:u-ef-inst-bound}
\end{equation}
\begin{equation}
U_{\mathrm{EE}}^{\mathrm{audit}} < p\cdot 2^{66}+p,
\label{eq:u-ee-inst-bound}
\end{equation}
\begin{equation}
U_{\mathrm{hor}}^{\mathrm{audit}} < p\cdot 2^{66}+p,
\label{eq:u-hor-inst-bound}
\end{equation}
\begin{equation}
U_{\mathrm{car}}^{\mathrm{audit}} < p\cdot 2^{31}+p.
\label{eq:u-car-inst-bound}
\end{equation}
\begin{equation}
U_{\mathrm{inv}}^{D,\mathrm{audit}} + p(2^{66}-1) < r,
\label{eq:u-inv-rhs-inst-bound}
\end{equation}
\begin{equation}
U_{\mathrm{EF}}^{D,\mathrm{audit}} + p(2^{31}-1) < r,
\label{eq:u-ef-rhs-inst-bound}
\end{equation}
\begin{equation}
U_{\mathrm{EE}}^{D,\mathrm{audit}} + p(2^{66}-1) < r,
\label{eq:u-ee-rhs-inst-bound}
\end{equation}
\begin{equation}
U_{\mathrm{hor}}^{D,\mathrm{audit}} + p(2^{66}-1) < r,
\label{eq:u-hor-rhs-inst-bound}
\end{equation}
\begin{equation}
U_{\mathrm{car}}^{D,\mathrm{audit}} + p(2^{31}-1) < r.
\label{eq:u-car-rhs-inst-bound}
\end{equation}
\end{lemma}

\begin{proof}
\eqref{eq:u-inv-inst-bound}, \eqref{eq:u-ef-inst-bound}, and
\eqref{eq:u-ee-inst-bound} are exactly the inverse/\(E\times F\)/\(E\times E\)
inequalities already verified in the concrete audit block above.

For \eqref{eq:u-hor-inst-bound}, by \eqref{eq:u-hor-inst} and
\eqref{eq:u-ee-inst},
\[
U_{\mathrm{hor}}^{\mathrm{audit}}
=
U_{\mathrm{EE}}^{\mathrm{audit}}
<
p\cdot 2^{66}+p.
\]
For \eqref{eq:u-car-inst-bound}, from \eqref{eq:u-car-inst} and
\[
p\cdot 2^{31}+p
=
4611686018427387903,
\]
we get
\[
4294967299 < 4611686018427387903.
\]
For RHS-side bounds, we use explicit instantiated values:
\[
U_{\mathrm{inv}}^{D,\mathrm{audit}} + p(2^{66}-1)
=
158456324954741698892249694207
<
2^{98},
\]
\[
U_{\mathrm{EF}}^{D,\mathrm{audit}} + p(2^{31}-1)
=
4611686016279904255
<
2^{98},
\]
\[
U_{\mathrm{EE}}^{D,\mathrm{audit}} + p(2^{66}-1)
=
158456324954741698892249694207
<
2^{98},
\]
\[
U_{\mathrm{hor}}^{D,\mathrm{audit}} + p(2^{66}-1)
=
316912649540548516720477732737
<
2^{98},
\]
\[
U_{\mathrm{car}}^{D,\mathrm{audit}} + p(2^{31}-1)
=
4611686018427387908
<
2^{98}.
\]
Since BN254 scalar \(r\) has 254 bits, \(2^{98}<r\). Hence
\eqref{eq:u-inv-rhs-inst-bound}--\eqref{eq:u-car-rhs-inst-bound} follow.
\end{proof}

\begin{corollary}[Instantiated no-wrap intervals for all audited families]
\label{cor:instantiated-no-wrap-intervals}
For every audited family
\(f\in\{f_{\mathrm{inv}},f_{\mathrm{EF}},f_{\mathrm{EE}},f_{\mathrm{hor}},f_{\mathrm{car}}\}\),
if a satisfying row obeys
\(\wt{L}_f\le U_f^{\mathrm{audit}}\),
\(\wt{D}_f\le U_f^{D,\mathrm{audit}}\), and \(\wt{Q}_f\in[0,B_f)\), then
\[
0\le \wt{L}_f < pB_f+p < r.
\]
\[
0\le \wt{D}_f + p\wt{Q}_f
\le
U_f^{D,\mathrm{audit}} + p(B_f-1)
< r.
\]
\end{corollary}

\begin{proof}
Immediate from Definition~\ref{def:instantiated-audit-constants},
Lemma~\ref{lem:instantiated-audit-closure}, and \(\wt{Q}_f\in[0,B_f)\).
\end{proof}

\begin{lemma}[Row-range discipline from audited templates]
\label{lem:row-range-from-audit}
Assume:
\begin{enumerate}[label=(\roman*),leftmargin=2em]
    \item Definitions~\ref{def:selector-relation-realization},
    \ref{def:template-realization-coherence}, and
    \ref{def:rowtable-zero-extension} hold, and every active local scalar
    relation instance is matched to one audited template \((f,\tau)\) with
    \(\tau\in\mathcal T_f\);
    \item all input wires satisfy the class/range constraints from
    Definitions~\ref{def:complete-quotient-coverage} and
    \ref{def:selector-residue-carry-wiring}, and
    \eqref{eq:residue-recompose}--\eqref{eq:u8-lookup};
    \item Definitions~\ref{def:template-audit-certificate} and
    \ref{def:instantiated-audit-constants} hold, and
    Lemma~\ref{lem:instantiated-audit-closure} applies.
\end{enumerate}
Then Definition~\ref{def:row-range-discipline} holds for all satisfying
witnesses. More precisely, for every row-table entry \((t,j)\), there exists
\(B_t\in\{2^{31},2^{66}\}\) such that
\[
0 \le \wt{L_t[j]} < pB_t + p < r.
\]
If \(j\in\Omega_t\), one may take \(B_t=B_f\) for the matched family \(f\); if
\(j\notin\Omega_t\), the claim is immediate from
Definition~\ref{def:rowtable-zero-extension}.
\end{lemma}

\begin{proof}
Fix any row-table entry \((t,j)\).

If \(j\notin\Omega_t\), then
Definition~\ref{def:rowtable-zero-extension} gives \(L_t[j]=0\), so
\[
0\le \wt{L_t[j]}=0 < pB_t+p<r
\]
for either choice \(B_t\in\{2^{31},2^{66}\}\).

Now assume \(j\in\Omega_t\). By (i), there is a unique active local scalar
relation instance realizing \((t,j)\), matched to some audited template
\((f,\tau)\). By Lemma~\ref{lem:exact-audited-template-evaluation},
\[
\wt{L_t[j]} = L_{f,\tau}(\wt{\mathbf z}_j)
\]
and
\[
0\le \wt{L_t[j]}\le U_f^{\mathrm{audit}}.
\]
By (iii) and Corollary~\ref{cor:instantiated-no-wrap-intervals},
\[
U_f^{\mathrm{audit}} < pB_f+p < r.
\]
Therefore
\[
0\le \wt{L_t[j]} < pB_f+p<r.
\]
Setting \(B_t:=B_f\) gives the claim.
\end{proof}

\begin{lemma}[RHS no-wrap from audited templates]
\label{lem:rhs-nowrap-from-audit}
Assume:
\begin{enumerate}[label=(\roman*),leftmargin=2em]
    \item all assumptions of Lemma~\ref{lem:row-range-from-audit};
    \item Definitions~\ref{def:complete-quotient-coverage},
    \ref{def:selector-quotient-wiring}, and
    \ref{def:family-quotient-class-coherence} hold;
    \item Lemma~\ref{lem:no-omitted-quotient-wire} applies.
\end{enumerate}
Then for each row relation \(L_t[j]=D_t[j]+pQ_t[j]\),
\[
0\le \wt{D_t[j]}+p\wt{Q_t[j]}<r.
\]
If \(j\in\Omega_t\) and the matched template family is \(f\), then one may use
that family's class \(B_f\); if \(j\notin\Omega_t\), the claim is immediate
from Definition~\ref{def:rowtable-zero-extension}.
\end{lemma}

\begin{proof}
Fix \((t,j)\).

If \(j\notin\Omega_t\), then Definition~\ref{def:rowtable-zero-extension}
gives \(D_t[j]=Q_t[j]=0\), so the claim is immediate.

Now assume \(j\in\Omega_t\) and let \(f\) be the family of the matched audited
template. By Lemma~\ref{lem:exact-audited-template-evaluation},
\[
0\le \wt{D_t[j]} = D_{f,\tau}(\wt{\mathbf z}_j)\le U_f^{D,\mathrm{audit}}.
\]
By Lemma~\ref{lem:no-omitted-quotient-wire}, \(Q_t[j]\) is range-constrained by
its slot class. By Definition~\ref{def:family-quotient-class-coherence}, that
slot class matches the family class \(B_f\). Hence
\[
\wt{Q_t[j]}\in[0,B_f).
\]
Therefore
\[
0\le \wt{D_t[j]}+p\wt{Q_t[j]}
\le
U_f^{D,\mathrm{audit}}+p(B_f-1).
\]
By Definition~\ref{def:template-audit-certificate} and
Lemma~\ref{lem:instantiated-audit-closure},
\[
U_f^{D,\mathrm{audit}}+p(B_f-1)<r.
\]
Combining gives the claim.
\end{proof}

\section{Coverage and wiring certificates}
\label{app:coverage-and-wiring}
This appendix records the explicit quotient inventory and selector-level
wiring certificates as public synthesis metadata.
\paragraph{Serialization granularity note.}
In the released certificate schema, wiring maps are serialized at
selector-slot granularity because the instantiated layout is row-invariant on
active rows. Accordingly, row indices are omitted from wiring-map blocks and
captured separately by \texttt{row\_activity} metadata.

\subsection{Coverage certificate (explicit inventory)}

We materialize Definition~\ref{def:complete-quotient-coverage} as a public
inventory certificate. The certificate lists every quotient/carry family,
its semantic origin, and its assigned range class.

\begin{center}
\begin{tabular}{|l|l|l|l|}
\hline
Family & Semantic origin & Class & Reason / bound target \\
\hline
\(Q_{\mathrm{inv}}\) & inverse rows \(L=q\cdot d\) & \(\mathcal Q_{66}\) & product can reach \(\Theta(p\cdot 2^{66})\) \\
\hline
\(Q_{\mathrm{EF},k}\) (\(k=0,\dots,3\)) & \(E\times F\) coordinate rows & \(\mathcal Q_{31}\) & \(L< p^2 < p\cdot 2^{31}+p\) \\
\hline
\(Q_{\mathrm{EE},k}\) (\(k=0,\dots,3\)) & \(E\times E\) coordinate rows & \(\mathcal Q_{66}\) & conservative high-range class \\
\hline
\(q_3,q_2,q_1\) & promoted Horner rows & \(\mathcal Q_{66}\) & promoted Horner subtemplate normalization quotients \\
\hline
\(Q_{\mathrm{car}}\) & carry-normalization rows & \(\mathcal Q_{31}\) & low-width carry transfer \\
\hline
\(c_2,c_1,v\) & boolean carry rows & \(\Cset_{\mathrm{bit}}\) & \(c(c-1)=0\) \\
\hline
\(u\) & 3-bit carry rows & \(\Cset_8\) & \(u\in\{0,\dots,7\}\) \\
\hline
\end{tabular}
\end{center}

\paragraph{Column-level class map (no omitted quotient symbol).}
The quotient inventory is explicitly decomposed as
\[
\mathcal Q
=
\mathcal Q_{\mathrm{inv}}
\sqcup
\mathcal Q_{\mathrm{EF}}
\sqcup
\mathcal Q_{\mathrm{EE}}
\sqcup
\mathcal Q_{\mathrm{hor}}
\sqcup
\mathcal Q_{\mathrm{car}},
\]
where
\begin{align*}
\mathcal Q_{\mathrm{inv}}
&:=
\{q_{\mathrm{inv}}^{(k)} \mid k\in\{0,1,2,3\}\},
\\
\mathcal Q_{\mathrm{EF}}
&:=
\{q_{\mathrm{EF}}^{(k)} \mid k\in\{0,1,2,3\}\},
\\
\mathcal Q_{\mathrm{EE}}
&:=
\{q_{\mathrm{EE}}^{(k)} \mid k\in\{0,1,2,3\}\},
\\
\mathcal Q_{\mathrm{hor}}
&:=
\{q_3^{(k)},q_2^{(k)},q_1^{(k)} \mid k\in\{0,1,2,3\}\},
\\
\mathcal Q_{\mathrm{car}}
&:=
\{q_{\mathrm{car}}^{(k)} \mid k\in\{0,1,2,3\}\}.
\end{align*}
Class assignments are fixed as
\[
\mathcal Q_{31}=\mathcal Q_{\mathrm{EF}}\sqcup\mathcal Q_{\mathrm{car}},
\qquad
\mathcal Q_{66}=\mathcal Q_{\mathrm{inv}}\sqcup\mathcal Q_{\mathrm{EE}}\sqcup\mathcal Q_{\mathrm{hor}}.
\]
Hence each promoted Horner row contributes exactly
\[
|\mathcal Q_{\mathrm{hor}}| = 12
\]
cells in \(\mathcal Q_{66}\), matching \(N_{66}^{\mathrm{round}}=12\) for the
Horner block. This scope excludes carry-normalization rows; their quotient
cells are \(q_{\mathrm{car}}^{(k)}\in\mathcal Q_{\mathrm{car}}\subset\mathcal Q_{31}\)
and are accounted separately in \(N_{31}\).
Any quotient symbol not in this certificate is invalid and circuit synthesis
rejects the instance.

\begin{definition}[Selector-level quotient wiring certificate]
\label{def:selector-quotient-wiring}
Let \(\Sigma\) be the set of enabled row selectors, and for each
\(s\in\Sigma\), let \(\Lambda_Q(s)\) be the local quotient slots used by the
row template of \(s\). Let \(\Omega_s\subseteq \Omega\) be the rows on which
selector \(s\) is active. The public metadata includes a deterministic map
\[
\omega_Q:
\{(s,\ell,j)\mid s\in\Sigma,\;\ell\in\Lambda_Q(s),\;j\in\Omega_s\}
\to
\mathcal Q.
\]
The certificate is valid only if:
\begin{enumerate}[label=(\roman*),leftmargin=2em]
    \item \textbf{Single-valued mapping:} each local quotient slot \((s,\ell,j)\) maps
    to exactly one global quotient column \(\omega_Q(s,\ell,j)\);
    \item \textbf{Class consistency:} each local slot has a fixed class tag
    \(b_{s,\ell}\in\{31,66\}\), and
    \[
    \kappa(\omega_Q(s,\ell,j)) = b_{s,\ell}
    \]
    for all active rows \(j\in\Omega_s\);
    \item \textbf{Exact coverage:}
    \[
    \operatorname{Im}(\omega_Q)=\mathcal Q;
    \]
    \item \textbf{No undeclared symbol:} synthesis rejects any quotient symbol
    not in \(\mathcal Q\).
\end{enumerate}
\end{definition}

\begin{lemma}[No omitted quotient wire]
\label{lem:no-omitted-quotient-wire}
Under Definitions~\ref{def:complete-quotient-coverage} and
\ref{def:selector-quotient-wiring}, every instantiated local quotient slot is
range-constrained by exactly one of
\eqref{eq:q31-recompose}--\eqref{eq:q31-lookups} or
\eqref{eq:q66-recompose}--\eqref{eq:q66-lookups}.
\end{lemma}

\begin{proof}
Fix an active local quotient slot \((s,\ell,j)\). By
Definition~\ref{def:selector-quotient-wiring}(i), it maps to a unique column
\(q:=\omega_Q(s,\ell,j)\in\mathcal Q\). By
Definition~\ref{def:complete-quotient-coverage}, \(q\) lies in exactly one of
\(\mathcal Q_{31}\) or \(\mathcal Q_{66}\), and each class is tied to exactly
one range-constraint system. Hence the slot is constrained by exactly one of
the two systems above. By (iii)--(iv), no instantiated quotient can bypass this
partition.
\end{proof}

\begin{definition}[Selector-level residue/carry wiring certificate]
\label{def:selector-residue-carry-wiring}
Let \(\Sigma\) be the set of enabled row selectors. For each
\(s\in\Sigma\), let \(\Lambda_N(s)\) be the local non-quotient slots of the row
template of \(s\) (residue or carry slots), and \(\Omega_s\subseteq\Omega\) the
rows where \(s\) is active. Each local slot \((s,\ell)\) has a fixed tag
\(\tau_{s,\ell}\in\{p,\mathrm{bit},8\}\). The public metadata includes a
deterministic map
\[
\omega_N:
\{(s,\ell,j)\mid s\in\Sigma,\;\ell\in\Lambda_N(s),\;j\in\Omega_s\}
\to
\Cset_p\sqcup\Cset_{\mathrm{bit}}\sqcup\Cset_8.
\]
The certificate is valid only if:
\begin{enumerate}[label=(\roman*),leftmargin=2em]
    \item \textbf{Single-valued mapping:} each local slot \((s,\ell,j)\) maps to exactly
    one global column \(\omega_N(s,\ell,j)\);
    \item \textbf{Tag consistency:}
    \[
    \tau_{s,\ell}=p \Rightarrow \omega_N(s,\ell,j)\in\Cset_p,\quad
    \tau_{s,\ell}=\mathrm{bit} \Rightarrow \omega_N(s,\ell,j)\in\Cset_{\mathrm{bit}},\quad
    \tau_{s,\ell}=8 \Rightarrow \omega_N(s,\ell,j)\in\Cset_8;
    \]
    \item \textbf{Exact coverage by tag:}
    \[
    \operatorname{Im}(\omega_N\!\mid_{\tau=p})=\Cset_p,\quad
    \operatorname{Im}(\omega_N\!\mid_{\tau=\mathrm{bit}})=\Cset_{\mathrm{bit}},\quad
    \operatorname{Im}(\omega_N\!\mid_{\tau=8})=\Cset_8;
    \]
    \item \textbf{No undeclared symbol:} synthesis rejects any residue/carry
    symbol outside \(\Cset_p\sqcup\Cset_{\mathrm{bit}}\sqcup\Cset_8\).
\end{enumerate}
\end{definition}

\begin{lemma}[No omitted residue/carry wire]
\label{lem:no-omitted-residue-carry-wire}
Under Definition~\ref{def:selector-residue-carry-wiring}, every instantiated
local residue/carry slot is range-constrained by exactly one of
\eqref{eq:residue-recompose}--\eqref{eq:residue-not-p},
\eqref{eq:bit-carry}, or \eqref{eq:u8-lookup}.
\end{lemma}

\begin{proof}
Fix any active local non-quotient slot \((s,\ell,j)\). By
Definition~\ref{def:selector-residue-carry-wiring}(i), it maps to a unique
column \(W:=\omega_N(s,\ell,j)\). By tag consistency (ii), \(W\) belongs to
exactly one of \(\Cset_p,\Cset_{\mathrm{bit}},\Cset_8\), and each class is tied
to exactly one binding rule family: residue constraints
\eqref{eq:residue-recompose}--\eqref{eq:residue-not-p}, bit carry
\eqref{eq:bit-carry}, or 3-bit carry lookup \eqref{eq:u8-lookup}. By (iii)--(iv),
no instantiated residue/carry slot can bypass this partition.
\end{proof}

\begin{definition}[Selector-level scalar-relation realization map]
\label{def:selector-relation-realization}
Let \(\Sigma\) be the set of enabled row selectors. For each \(s\in\Sigma\),
let \(\Lambda_R(s)\) be the set of local scalar relation slots in selector
template \(s\), and let \(\Omega_s\subseteq\Omega\) be the rows where \(s\) is
active. Let \(\mathcal I_{\mathrm{rel}}\) be the global relation-index set for
row tables \((L_t[j],D_t[j],Q_t[j])_{(t,j)\in\mathcal I_{\mathrm{rel}}\times\Omega}\).

The public metadata includes a deterministic map
\[
\omega_R:
\{(s,\rho,j)\mid s\in\Sigma,\;\rho\in\Lambda_R(s),\;j\in\Omega_s\}
\to
\mathcal I_{\mathrm{rel}}.
\]
The certificate is valid only if:
\begin{enumerate}[label=(\roman*),leftmargin=2em]
    \item \textbf{Single-valued mapping:} each active local scalar relation instance
    \((s,\rho,j)\) maps to exactly one relation index
    \(t=\omega_R(s,\rho,j)\);
    \item \textbf{Exact active coverage / no duplication:} for each
    \((t,j)\) with \(j\in\Omega_t\), there exists a unique active local scalar
    relation instance \((s,\rho,j)\) such that \(\omega_R(s,\rho,j)=t\);
    \item \textbf{Realization rule:} for that unique \((s,\rho,j)\), the row
    table values satisfy
    \[
    L_t[j]=\mathcal L_{s,\rho}(\mathbf z_j),\quad
    D_t[j]=\mathcal D_{s,\rho}(\mathbf z_j),\quad
    Q_t[j]=\mathcal Q_{s,\rho}(\mathbf z_j),
    \]
    where \(\mathbf z_j\) is the row-local witness tuple and
    \(\mathcal L_{s,\rho},\mathcal D_{s,\rho},\mathcal Q_{s,\rho}\) are the
    declared local realization expressions for slot \((s,\rho)\);
    \item \textbf{No undeclared relation:} synthesis rejects any active local
    scalar relation instance not covered by \(\omega_R\).
\end{enumerate}
\end{definition}

\begin{definition}[Template/realization coherence]
\label{def:template-realization-coherence}
Let \(\Sigma\) be the enabled selector set from
Definition~\ref{def:selector-relation-realization}. The public metadata
includes a deterministic map
\[
\omega_T:
\{(s,\rho,j)\mid s\in\Sigma,\;\rho\in\Lambda_R(s),\;j\in\Omega_s\}
\to
\{(f,\tau)\mid f\in\mathcal F,\;\tau\in\mathcal T_f\}.
\]
For each active local scalar relation instance \((s,\rho,j)\), define
\[
t:=\omega_R(s,\rho,j),
\qquad
(f,\tau):=\omega_T(s,\rho,j).
\]
Let \(\mathbf z_j\) be the row-local field tuple used by the local realization
expressions, and let \(\wt{\mathbf z}_j\) be the tuple of canonical integer
lifts of its ranged entries. Coherence requires
\[
L_t[j]
=
\mathcal L_{s,\rho}(\mathbf z_j)
=
\bigl[L_{f,\tau}(\wt{\mathbf z}_j)\bigr]_R,
\]
\[
D_t[j]
=
\mathcal D_{s,\rho}(\mathbf z_j)
=
\bigl[D_{f,\tau}(\wt{\mathbf z}_j)\bigr]_R,
\]
where \([\cdot]_R\) denotes reduction modulo \(r\) followed by the canonical
embedding into \(R\). Quotient slots appearing in the same instantiated
template must satisfy Definition~\ref{def:family-quotient-class-coherence}.
\end{definition}

\paragraph{Worked end-to-end example (an \(E\times E\) coordinate row).}
Fix an active row \(j^\star\in\Omega\) and a coordinate \(k\in\{0,1,2,3\}\).
Consider a local scalar relation instance \((s,\rho,j^\star)\) whose selector
slot corresponds to the \(E\times E\) coordinate relation for coordinate \(k\),
and let
\[
t:=\omega_R(s,\rho,j^\star),
\qquad
(f,\tau):=\omega_T(s,\rho,j^\star).
\]
In this case we have \(f=\mathrm{EE}\) and one can take \(\tau=k\). Let the
row-local witness tuple be
\[
\mathbf z_{j^\star}=(x_0,\dots,x_3,\;y_0,\dots,y_3,\;c)\in R^{9},
\]
where canonical binding implies there exist integer lifts
\(\wt{x}_a,\wt{y}_b,\wt{c}\in I_p\) for all \(a,b\in\{0,1,2,3\}\). The declared
local realization expressions for this slot take the concrete form
\[
L_t[j^\star]
=
\mathcal L_{s,\rho}(\mathbf z_{j^\star})
=
\sum_{a,b=0}^{3} m_{ab}^k\,x_a y_b,
\qquad
D_t[j^\star]
=
\mathcal D_{s,\rho}(\mathbf z_{j^\star})
=
c,
\]
with quotient value \(Q_t[j^\star]=\mathcal Q_{s,\rho}(\mathbf z_{j^\star})\)
wired to the family slot \(Q_{\mathrm{EE},k}\). Template/realization coherence
then means these match the audited normalized templates
\[
L_{\mathrm{EE},k}(\wt{\mathbf z}_{j^\star})
:=
\sum_{a,b=0}^{3} m_{ab}^k\,\wt{x}_a \wt{y}_b,
\qquad
D_{\mathrm{EE},k}(\wt{\mathbf z}_{j^\star})
:=
\wt{c},
\]
via
\(
L_t[j^\star]=[L_{\mathrm{EE},k}(\wt{\mathbf z}_{j^\star})]_R
\)
and
\(
D_t[j^\star]=[D_{\mathrm{EE},k}(\wt{\mathbf z}_{j^\star})]_R
\).
The audit certificate bounds the coefficients by
\(0\le m_{ab}^k\le M_{\mathrm{EE}}^{\max}\le p-1\) and assigns the family class
\(B_{\mathrm{EE}}=2^{66}\), yielding the certified family bounds
\(
U_{\mathrm{EE}}^{\mathrm{audit}}=16(p-1)^3
\)
and
\(
U_{\mathrm{EE}}^{D,\mathrm{audit}}=p-1
\).
Here \(m_{ab}^k\) denotes the audited non-negative template coefficient in the
instantiated library, not a raw signed basis coefficient.
Finally, the coverage/wiring certificate assigns \(Q_{\mathrm{EE},k}\) to a
\(\mathcal Q_{66}\) column, so canonical binding gives
\(\wt{Q_t[j^\star]}\in[0,2^{66})\). This is the concrete
``selector \(\to\) realization \(\to\) normalized \((L,D,Q)\) \(\to\) bound
\(\to\) quotient class'' chain used by the exactness argument.

\begin{lemma}[Exact evaluation of audited templates on active rows]
\label{lem:exact-audited-template-evaluation}
Assume Definitions~\ref{def:selector-relation-realization},
\ref{def:template-realization-coherence}, and
\ref{def:template-audit-certificate}. Let \((t,j)\) be an active row-table
entry, let \((s,\rho,j)\) be its unique realizing local scalar relation
instance, and let \((f,\tau)\) be the unique audited template assigned to that
instance. If all row-local inputs lie in their certified integer intervals, then
\[
L_t[j]
=
\bigl[L_{f,\tau}(\wt{\mathbf z}_j)\bigr]_R,
\qquad
D_t[j]
=
\bigl[D_{f,\tau}(\wt{\mathbf z}_j)\bigr]_R,
\]
with
\[
0\le L_{f,\tau}(\wt{\mathbf z}_j)\le U_f^{\mathrm{audit}},
\qquad
0\le D_{f,\tau}(\wt{\mathbf z}_j)\le U_f^{D,\mathrm{audit}}.
\]
Consequently, whenever
\[
U_f^{\mathrm{audit}}<r,
\qquad
U_f^{D,\mathrm{audit}}<r,
\]
the field values \(L_t[j]\) and \(D_t[j]\) are the images in \(R\) of those
exact integers, i.e.
\[
\wt{L_t[j]} = L_{f,\tau}(\wt{\mathbf z}_j),
\qquad
\wt{D_t[j]} = D_{f,\tau}(\wt{\mathbf z}_j).
\]
\end{lemma}

\begin{proof}
The displayed field identities are exactly the content of
Definition~\ref{def:template-realization-coherence}. By
Lemma~\ref{lem:template-bound-soundness},
\[
0\le L_{f,\tau}(\wt{\mathbf z}_j)\le U_f^{\mathrm{audit}},
\qquad
0\le D_{f,\tau}(\wt{\mathbf z}_j)\le U_f^{D,\mathrm{audit}}.
\]
Expand \(L_{f,\tau}\) and \(D_{f,\tau}\) as in
Definition~\ref{def:template-audit-certificate}: they are sums of linear terms
\(a\,z_u\) and quadratic terms \(b\,z_u z_v\) with non-negative integer
coefficients. Since \(U_f^{\mathrm{audit}}<r\) and \(U_f^{D,\mathrm{audit}}<r\),
also the per-template bounds \(U_{f,\tau}^{L}\) and \(U_{f,\tau}^{D}\) are
\(<r\). By Definition~\ref{def:template-audit-certificate}, every linear
contribution \(a_{f,\tau,u}U_{f,u}^{\max}\) and quadratic contribution
\(b_{f,\tau,uv}U_{f,u}^{\max}U_{f,v}^{\max}\) is a non-negative summand of
\(U_{f,\tau}^{L}\) or \(U_{f,\tau}^{D}\), hence is \(<r\). Lemma~\ref{lem:monomial-nowrap}
therefore applies to each multiplication. For every nonzero quadratic term,
\(b_{f,\tau,uv}\ge 1\), so
\[
U_{f,u}^{\max}U_{f,v}^{\max}
\le
b_{f,\tau,uv}U_{f,u}^{\max}U_{f,v}^{\max}
<
r.
\]
Hence Lemma~\ref{lem:monomial-nowrap} applies first to \(z_u z_v\), and then
again to multiplication by \(b_{f,\tau,uv}\). Zero-coefficient terms are
omitted from the template representation. Thus every summand evaluates in \(R\)
as the image of its exact integer value. Because all summands are non-negative
and the total is \(<r\), addition does not wrap either, so the field evaluations
coincide with the images in \(R\) of the total integer evaluations. The
canonical lifts are therefore exactly the displayed integers.
\end{proof}

\begin{lemma}[No omitted active local relation]
\label{lem:no-omitted-active-relation}
Under Definition~\ref{def:selector-relation-realization}, every active local
scalar relation instance contributes to exactly one row-table entry
\((L_t[j],D_t[j],Q_t[j])\), and every active row-table entry arises from exactly
one such instance.
\end{lemma}

\begin{proof}
Immediate from
Definition~\ref{def:selector-relation-realization}(i)--(ii).
\end{proof}

\begin{definition}[Row-table activity partition and zero extension]
\label{def:rowtable-zero-extension}
For each relation index \(t\), there is a public active-row set
\(\Omega_t\subseteq\Omega\), and row-table columns
\((L_t[j],D_t[j],Q_t[j])_{j\in\Omega}\). The circuit enforces:
\begin{enumerate}[label=(\roman*),leftmargin=2em]
    \item \textbf{Active-row realization:} for each \(j\in\Omega_t\), the tuple
    \((L_t[j],D_t[j],Q_t[j])\) is instantiated from exactly one active local
    scalar relation instance, consistent with
    Definition~\ref{def:selector-relation-realization};
    \item \textbf{Inactive-row fixing:} for each \(j\in\Omega\setminus\Omega_t\),
    \[
    L_t[j]=0,\qquad D_t[j]=0,\qquad Q_t[j]=0.
    \]
\end{enumerate}
\end{definition}

\begin{lemma}[No inactive-row influence on compressed evaluation]
\label{lem:no-inactive-row-influence}
Under Definitions~\ref{def:rowtable-zero-extension} and
\ref{def:compression-rowwise-rhs-consistency}, for each relation index \(t\) and
each \(j\in\Omega\setminus\Omega_t\),
\[
L_t(\omega^j)-D_t(\omega^j)-pQ_t(\omega^j)=0.
\]
Consequently, any non-zero residual entry occurs only on active rows.
\end{lemma}

\begin{proof}
Fix \(t\) and \(j\in\Omega\setminus\Omega_t\). By
Definition~\ref{def:rowtable-zero-extension}, \(L_t[j]=D_t[j]=Q_t[j]=0\). By
Definition~\ref{def:compression-rowwise-rhs-consistency},
\(L_t(\omega^j)=L_t[j]\), \(D_t(\omega^j)=D_t[j]\), and \(Q_t(\omega^j)=Q_t[j]\).
Substituting gives the claimed zero residual.
\end{proof}

\section{Supplementary implemented B1-lite circuit-shape comparison}
\label{app:b1lite-supplement}
This appendix records the implemented Halo2 harness comparison referenced in
the main text. It remains a circuit-shape comparison under one backend, not a
task-equated benchmark.

\begin{table}[t]
\centering
\small
\setlength{\tabcolsep}{6pt}
\begin{tabular}{|l|r|r|r|}
\hline
Metric & Repair fragment & Repair B2 paired package & B1-lite baseline \\
\hline
Rows & 72 & 76 & 352 \\
Lookup cells & 152 & 160 & 126 \\
Native multiplicative constraints & 56 & 56 & 256 \\
Linear constraints & 56 & 60 & 63 \\
\hline
\end{tabular}
\caption{Implemented Halo2 circuit-shape comparison from the artifact harness.
The B1-lite baseline includes explicit \(16\times16\) limb cross products and
carry-chain constraints for a \(256\)-bit multiplication, but excludes modular
reduction and CRT-consistency checks; thus it is a conservative partial
baseline for the implemented schoolbook-limb wrong-field family.}
\label{tab:b1lite-implemented-comparison-appendix}
\end{table}

In this harness, B1-lite versus the repair fragment gives
\(352/72=4.89\times\) in rows and \(256/56=4.57\times\) in multiplicative
constraints; versus the Repair B2 paired package, it gives
\(352/76=4.63\times\) in rows and again \(256/56=4.57\times\) in multiplicative
constraints. The row counts are backend-specific physical-row counts and are
not expected to be a simple function of symbolic arithmetic-object counts,
because lookup batching and gate packing can place multiple objects in one row.
The lookup-cell ratio goes in the opposite direction, matching the design goal
of this repair layer: trading multiplicative constraints for batched native
lookups.\footnote{Artifact repository:
\url{https://github.com/junbyjun1238/zk2-bench-b1-lite}.
Relevant files are \texttt{src/main.rs} and \texttt{results.md}. Reproduce by
running \texttt{cargo run --release}. Exact artifact binding for this note is
given by the published certificate/manifests and checker command in
Appendix~\ref{app:audit-certificate}.}

% END OPTION2 APPENDICES

\begin{thebibliography}{99}

\bibitem{plonk2019}
Ariel Gabizon, Zachary J. Williamson, and Oana Ciobotaru.
\newblock PLONK: Permutations over Lagrange-bases for Oecumenical Noninteractive arguments of Knowledge.
\newblock IACR ePrint 2019/953, 2019.
\newblock \url{https://eprint.iacr.org/2019/953}.

\bibitem{plookup2020}
Ariel Gabizon and Zachary J. Williamson.
\newblock Plookup: A simplified polynomial protocol for lookup tables.
\newblock IACR ePrint 2020/315, 2020.
\newblock \url{https://eprint.iacr.org/2020/315}.

\bibitem{lookup-sok-2025}
Hossein Hafezi, Gaspard Anthoine, Matteo Campanelli, and Dario Fiore.
\newblock SoK: Lookup Table Arguments.
\newblock IACR ePrint 2025/1876, 2025.
\newblock \url{https://eprint.iacr.org/2025/1876}.

\bibitem{foreignfield2025}
Miguel Ambrona, Denis Firsov, and Inigo Querejeta-Azurmendi.
\newblock Efficient Foreign-Field Arithmetic in PLONK.
\newblock IACR ePrint 2025/695, 2025.
\newblock \url{https://eprint.iacr.org/2025/695}.

\bibitem{halo2wrong}
Privacy Scaling Explorations.
\newblock halo2wrong: wrong-circuit gadgets for Halo2 (including non-native arithmetic).
\newblock GitHub repository, commit
\texttt{3eaa2c4703a333e9695dc980bebf3445b7a9b141}.
\newblock \url{https://github.com/privacy-scaling-explorations/halo2wrong/tree/3eaa2c4703a333e9695dc980bebf3445b7a9b141}.

\end{thebibliography}
\end{document}
